\documentclass[12pt, twoside]{article}
% Emit local LDF shims before packages load
\input{lib/datatool-ldf}

% Make it possible to conditionally depend on the TeX engine used
\RequirePackage{ifluatex}

\RequirePackage{comment}

% To access the \captionof{} functionality
\RequirePackage{caption}

%\RequirePackage{silence}
%\WarningFilter{latexfont}{Font shape `U/stmry/m/n' in size <5.5> not available}
%\WarningFilter{latexfont}{Font shape}

\newif\ifinswedish
\inswedishfalse

\RequirePackage{etoolbox}

% The following toggles are needed because of kth/kthcolors.tex and lib/includes.tex
\newif\ifnomenclature
\nomenclaturefalse

\newif\ifdigitaloutput
\digitaloutputfalse

\makeatletter
%% Define a pair of commands to disable and reenable specific packages - see https://tex.stackexchange.com/questions/39415/unload-a-latex-package
\makeatletter
\newcommand{\disablepackage}[2]{%
  \disable@package@load{#1}{#2}%
}
\newcommand{\reenablepackage}[1]{%
  \reenable@package@load{#1}%
}
\makeatother

%% To avoid the warning: "Package transparent Warning: Loading aborted, because pdfTeX is not running in PDF mode."
\disablepackage{transparent}{}

% include a variety of packages that are useful
\input{lib/includes}

\usepackage{minted}
\RequirePackage{fancyhdr} % Take control of headers and footers

%\usepackage{luatexja-fontspec}
%\ltjdefcharrange{8}{"2190-"21FF} % Define range 8 as Arrows
%\ltjsetparameter{jacharrange={-8}} % Treat range 8 as Western characters (ALChar)

%\ltjsetparameter{jacharrange={-3}} % Treat range 3 as Western characters (ALChar)

\begin{comment}
\ifinswedish
    \usepackage[english, main=swedish, bidi=basic]{babel}
\else
    \usepackage[swedish, main=english, provide+=*, bidi=basic]{babel}
\fi
\end{comment}
\ifinswedish
    \usepackage[main=swedish]{babel}
\else
    \usepackage[main=english, provide+=*]{babel}
\fi
%\setmainjfont{Noto Sans CJK JP}

%\usepackage[style=ieee,citestyle=numeric-comp, maxbibnames=99, giveninits=false]{biblatex}
% Input centralized biblatex configuration
\input{lib/biblatex-and-friends}



% Use the chapter* style for the heading of the bibliography
\defbibheading{bibliography}[\bibname]{%
\section*{#1}%   skip \centering 
}
  
\addbibresource{references.bib}    
\input{kth/kthcolors}



\input{lib/defines}  % load some additional definitions to make writing more consistent

\usepackage[plainpages=false, unicode=true]{hyperref}
\input{lib/includes-after-hyperref}
\usepackage{orcidlink}  %% to provide ORCID icons and links

%\usepackage[acronym, style=super, toc=false, nonumberlist, nomain, nopostdot=true, notranslate]{glossaries-extra}
% In the document preamble where glossaries-extra is loaded:
\usepackage[
  acronym,
  toc=false,
  nonumberlist,
  nomain,
  nopostdot=true, 
  nohypertypes={main, acronym}, % disables link generation to targets that do not exist
  notranslate,
]{glossaries-extra}
%% In overleaf, if this file is in a folder and not the root
%% using \makeglossaries will not work, as the Overleaf latexmk
%% will not run the program to take the output.acn file and make the output.acr file
%\makeglossaries
% However, you can get TeX to do the work with the following:
\makenoidxglossaries
% For details of the Overleaf make process, see https://www.overleaf.com/learn/how-to/How_does_Overleaf_compile_my_project%3F
%% Potentially, one could set up a latexmkrc file in the folder

\input{lib/acronyms}                %load the acronyms file

% Include Glossary ---
% Align the text expansion of the glossary entries
\newglossarystyle{mylong}{%
  \setglossarystyle{long}%
  \renewenvironment{theglossary}%
     {\begin{longtable}[l]{@{}p{\dimexpr 2cm-\tabcolsep}p{0.8\hsize}}}% <-- change the value here
     {\end{longtable}}%
 }

\input{lib/placeHolder_lbx_files}

% define a left-aligned table cell that is ragged right
\newcolumntype{L}[1]{>{\raggedright\let\newline\\\arraybackslash\hspace{0pt}}p{#1}}

\usepackage{cleveref}           %% Replace Section with a symbol
\usepackage{metalogo}   % for \LuaLaTeX logo

  \RequirePackage{fontspec}
  \defaultfontfeatures{Ligatures={TeX}} % This enables TeX style ligatures such as ---, '', ``, and so on




    \babelfont{rm}{TeX Gyre Termes}
    \DeclareFontShape{TU}{TeXGyreTermes(0)}{md}{n}{<->sub * TeXGyreTermes(0)/m/n}{}
    \DeclareFontShape{TU}{TeXGyreTermes(0)}{sb}{n}{<->ssub * TeXGyreTermes(0)/b/n}{}
    \DeclareFontShape{TU}{TeXGyreTermes(0)}{sb}{it}{<->ssub * TeXGyreTermes(0)/b/it}{}
  
    %\setsansfont{TeX Gyre Heros}   %% Helvetica like font
    \babelfont{sf}{TeX Gyre Heros}
    %\setmonofont[Ligatures={NoCommon}, Numbers={Lining,Monospaced}]{TeX Gyre Cursor}  %% Courier like font
    %% Turn off ligature for tt font
    \babelfont{tt}[Ligatures={NoCommon, NoRequired, NoContextual}]{TeX Gyre Cursor}
 \begin{comment}
    \babelfont[english]{rm}{TeX Gyre Termes}
    \babelfont[english]{sf}{TeX Gyre Heros}
    \babelfont[english]{tt}{TeX Gyre Cursor}
\end{comment}

 %  \setmathfont{TeX Gyre Termes Math} %% a math font
    \usepackage{mathtools}
 \usepackage[warnings-off={mathtools-colon,mathtools-overbracket}]{unicode-math}
  % The [version=bold, FakeBold=1.2] is to avoid a warning about the lack of a bold font
  %\setmathfont{TeX Gyre Pagella Math}[version=bold, FakeBold=1.2] %% a font for math
  \setmathfont{STIX Two Math}[version=normal]
  %\setmathfont{TeX Gyre Pagella Math}[version=normal]
  %\setmathfont{TeX Gyre Pagella Math}[version=bold, BoldFeatures={FakeBold=1.5}]
  % For both XeLaTeX and LuaLatex for getting access to unicode symbols
  %\newfontfamily\myfont[CharacterVariant=1]{NewCM10-Regular.otf}
  % STIX Project (Scientific and Technical Information Exchange)
  % STIX Two Math does not have bold face - so we fake it
  %\newfontfamily\mystixmathfont[BoldFeatures={FakeBold=1.5}, BoldItalicFeatures={FakeBold=1.5}]{STIX Two Math}
  \newfontfamily\mystixmathfont{STIX Two Math}
  % STIX Two Math does not have a bold font, but it has bold symbols with and without serifs - but you manually have to use them, unless you are in math mode - then you can use \symbf{}
  % and this will return the bold serif version of the character
\makeatletter
\AtBeginDocument{%
  \@for\idx:={0,1,2,3,4,5,6}\do{%
    \@ifundefined{T@TU+STIXTwoMath(\idx)}{%
      % Family not loaded by unicode-math; do nothing
    }{%
      \DeclareFontShape{TU}{STIXTwoMath(\idx)}{b}{n}{<->ssub * STIXTwoMath(\idx)/m/n}{}%
      \DeclareFontShape{TU}{STIXTwoMath(\idx)}{sb}{n}{<->ssub * STIXTwoMath(\idx)/m/n}{}%
      \DeclareFontShape{TU}{STIXTwoMath(\idx)}{b}{it}{<->ssub * STIXTwoMath(\idx)/m/it}{}%
      \DeclareFontShape{TU}{STIXTwoMath(\idx)}{sb}{it}{<->ssub * STIXTwoMath(\idx)/m/it}{}%
    }%
  }%
}
\makeatother
  \newfontfamily\mystixtextfont{STIX Two Text}
  
    % To access the Stylistic Set 1 <ss01> such as ℋ
    \newfontfamily\mystixmathfontSSa{STIX Two Math}[StylisticSet=1]

    % use english as a fallback when in other languages
    \babelprovide[import, onchar=ids fonts]{english}

    
  % for new KTH cover
  % Load the Figtree font as it is used for the new KTH graphical profile
  % 

    \newfontfamily{\FigtreeFont}[Ligatures=TeX,
        Path=./Figtree/static/,
        Extension = .ttf,
        UprightFont=*-Regular,
        BoldFont=*-Bold,
        BoldItalicFont=*-BoldItalic,
        ItalicFont=*-Italic,
        %FontFace={l}{n}{*-Light},
        %FontFace={l}{it}{*-LightItalic},
        FontFace={md}{n}{*-Medium},
        FontFace={md}{it}{*-MediumItalic},
        FontFace={sb}{n}{*-Semibold},
        FontFace={sb}{it}{*-SemiBoldItalic},
        %FontFace={k}{n}{*-Black},
        %FontFace={k}{it}{*-BlackItalic},
        %FontFace={eb}{n}{Font=*-ExtraBold},
        %FontFace={eb}{it}{Font=*-ExtraBoldItalic}
        ]{Figtree}


\newfontfamily\pageNumberFont{Figtree} %% set the font to use for page numbering

\newfontfamily{\NotoEmojiFont}[Ligatures=TeX,
    Path=./Noto_Emoji/static/,
    Extension = .ttf,
    UprightFont=*-Regular,
    BoldFont=*-Bold,
    FontFace={l}{n}{*-Light.ttf},
    FontFace={md}{n}{*-Medium},
    FontFace={sb}{n}{*-SemiBold},
    ]{NotoEmoji}



\begin{comment}
   % To set the abstract headings in Figtree, we redefine the abstract environment to look at the language being used and use the appropriate font, with the default being Figtree
   % The languages that are automatically introduced by Polyglossia or Babel have a name of the form xxxfont and xxxfontsf; where xxxfont is the serif font and xxxfontsf is the sans serif font.
   % This means that for each language that Figtree does not support, you have to define the sans serif and serif font to use.

      \babelprovide[import, onchar=ids fonts]{greek}
      \babelprovide[import, onchar=ids fonts]{russian}
      \babelprovide[import, onchar=ids fonts]{vietnamese}
      \babelfont[greek, russian, vietnamese]{rm}{Noto Serif}
      \babelfont[greek, russian, vietnamese]{sf}{Noto Sans}
      \babelfont[greek, russian, vietnamese]{tt}{Noto Mono}

  
    \babelprovide[import, onchar=ids fonts]{hindi}
    \babelfont[hindi]{rm}{Noto Serif Devanagari}
    \babelfont[hindi]{sf}{Noto Sans Devanagari}
    \babelfont[hindi]{tt}{Noto Sans Devanagari} % Noto Mono does not have the glyphs

    %\babelprovide[import, onchar=ids fonts]{russian}
    %\babelfont[russian]{rm}{Noto Serif}
    %\babelfont[russian]{sf}{Noto Sans}
    %\babelfont[russian]{tt}{Noto Mono}

    \babelprovide[import, onchar=ids fonts]{chinese-simplified}
    \babelfont[chinese-simplified]{rm}{Noto Serif CJK SC}
    \babelfont[chinese-simplified]{sf}{Noto Sans CJK SC}
    \babelfont[chinese-simplified]{tt}{Noto Sans Mono CJK SC}
    
    \babelfont[chinese-traditional]{rm}{Noto Serif CJK TC}
    \babelfont[chinese-traditional]{sf}{Noto Sans CJK TC}
    \babelfont[chinese-traditional]{tt}{Noto Sans Mono CJK TC}
  
    \babelprovide[import, onchar=ids fonts]{japanese}
    \babelfont[japanese]{rm}{Noto Serif CJK JP}
    \babelfont[japanese]{sf}{Noto Sans CJK JP}
    \babelfont[japanese]{tt}{Noto Sans Mono CJK JP}
  
  % If you are going to use Arabic

    \babelprovide[import, onchar=ids fonts]{arabic}
    \babelprovide[import, onchar=ids fonts]{centralkurdish}
    \babelfont[arabic, centralkurdish]{rm}{Noto Naskh Arabic}
    \babelfont[arabic, centralkurdish]{sf}{Noto Sans Arabic}
    \babelfont[arabic, centralkurdish]{tt}{Noto Sans Arabic}
    % If one really needs a monospaced font, one might try Kawkab Mono
    % However, it seems that it is a work in progress - see https://makkuk.com/kawkab-mono/ and https://github.com/aiaf/kawkab-mono/tree/master

    %\babelprovide[import, onchar=ids fonts]{centralkurdish}
    %\babelfont[centralkurdish]{rm}{Noto Naskh Arabic}
    %\babelfont[centralkurdish]{sf}{Noto Sans Arabic}
    %\babelfont[centralkurdish]{tt}{Noto Sans Arabic}

      
  % If you are going to use Hebrew or Yiddish
\babelprovide[import, onchar=ids fonts]{hebrew}
\babelprovide[import, onchar=ids fonts]{yiddish}
\babelfont[hebrew, yiddish]{rm}{Noto Serif Hebrew}
\babelfont[hebrew, yiddish]{sf}{Noto Sans Hebrew}
\babelfont[hebrew, yiddish]{tt}{Noto Sans Hebrew}


    %\babelprovide[import, onchar=ids fonts]{vietnamese}
    %\babelfont[vietnamese]{rm}{Noto Serif}
    %\babelfont[vietnamese]{sf}{Noto Sans}
    %\babelfont[vietnamese]{tt}{Noto Mono}

\end{comment}
  
    % The Overleaf TeX Live includes these fonts, so there is little you have to do!
    % The list of such fonts is at https://www.overleaf.com/learn/latex/Questions%2FWhich_OTF_or_TTF_fonts_are_supported_via_fontspec%3F
    %
\newfontfamily{\NotoSansJPMonoFont}[Ligatures=TeX,
    ]{Noto Sans Mono CJK JP Regular}


\newfontfamily{\NotoSansJPFont}[Ligatures=TeX,
    ]{Noto Sans CJK JP}


\newfontfamily{\NotoSansFont}[Ligatures=TeX,
Extension = .ttf,
    UprightFont = NotoSans-Regular,
    BoldFont = NotoSans-Bold,
    ItalicFont = NotoSans-Italic,
    BoldItalicFont = NotoSans-BoldItalic
    ]{NotoSansCustomA}
    
\newfontfamily{\NotoSerifFont}[Ligatures=TeX,
Extension = .ttf,
    UprightFont = NotoSerif-Regular,
    BoldFont = NotoSerif-Bold,
    ItalicFont = NotoSerif-Italic,
    BoldItalicFont = NotoSerif-BoldItalic
    ]{NotoSerifCustomA}

\newfontfamily{\DejaVuSansFont}[Ligatures=TeX,]{DejaVu Sans}

% A font to be used in minted listings
\newfontfamily{\ListingMono}{Noto Sans Mono}[Ligatures=ResetAll]

% fonts for various symbols
\newfontfamily{\NotoSansSymbolsFont}{Noto Sans Symbols}[Ligatures=ResetAll]
\newfontfamily{\NotoSansSymbolsTwoFont}{Noto Sans Symbols 2}[Ligatures=ResetAll]



% color symbols - this would require changes in the font used in the U+2600-U+26FF-Miscellaneous-Symbols.tex file
% for example for characters: ☕, ⚽, ⛄, and ✈
\newfontfamily{\NotoColorEmojiFont}{Noto Color Emoji}[Renderer=Harfbuzz]

% Phonetic Extensions, U+1D00 - U+1D7F
\input{unicode_blocks/U+1D00-U+1D7F-Phonetic-Extensions}
% Superscripts and Subscripts, U+2070 - U+209F
\input{unicode_blocks/U+2070-U+209F-Superscripts-and-Subscripts}
% Miscellaneous Symbols, U+2600 - U+26FF
\input{unicode_blocks/U+2600-U+26FF-Miscellaneous-Symbols}
% Dingbats, U+2700 - U+27BF
\input{unicode_blocks/U+2700-U+27BF-Dingbats}
% Miscellaneous Symbols and Pictographs, U+1F300 - U+1F5FF
\input{unicode_blocks/U+1F300-U+1F5FF-Miscellaneous-Symbols-and-Pictographs}

% Set up the header and footer for plain style pages
\fancypagestyle{plain}{%
\fancyhf{}% clear all header and footer fields
\fancyfoot[C]{\pageNumberFont\small\selectfont\thepage}
\renewcommand{\headrulewidth}{0pt}%
\renewcommand{\footrulewidth}{0pt}%
}
% Set up the header and footer
\fancyhead{}
%\fancyhead[RO]{\sffamily\small\leftmark\thinspace|\thinspace\thepage}
%\fancyhead[LE]{\sffamily\small\thepage\thinspace|\thinspace\leftmark}
% Without conversion to uppercase
% Note that we need to switch to the familydefault for the page numbers
\fancyhead[RO]{{\sffamily\small\nouppercase\leftmark}{\pageNumberFont\small\selectfont \thinspace|\thinspace \thepage}}
\fancyhead[LE]{{\pageNumberFont\small\selectfont\thepage \thinspace|\thinspace}{\sffamily\small\nouppercase\leftmark}}
\fancyfoot{}
\renewcommand{\headrulewidth}{0pt}
\pagestyle{fancy}
\renewcommand{\sectionmark}[1]{%
\markboth{#1}{}}


\setlength{\headheight}{15pt}
\addtolength{\topmargin}{-3pt}

\input{lib/configure_listings}
\sbox{\lstbreakbox}{\textcolor{red}{\ensuremath{\mystixmathfont ↪}}\space}
% If you use both listing and lstlisting environments - the following makes them both use the same counter and the same formatting in the List of Listings
\makeatletter
\AtBeginDocument{\let\c@listing\c@lstlisting}
\AtBeginDocument{\let\l@listing\l@lstlisting}
\makeatother

\newcommand{\dname}[1]{\textbf{#1}}
\newcommand{\fname}[1]{\texttt{#1}}


\RequirePackage{luacode}
\input{kth/Lua_timestamping}

% Input file generators before \documentclass
\input{lib/ext-eprint-dbx}

\input{lib/check_for_placeholder_references}

\title{README and notes about the template for programmers -- Biber and \BibLaTeX{}}

\author{Gerald Q. Maguire Jr.}
\date{October 2026}

\begin{document}
\pagenumbering{roman}
\glsresetall
\maketitle
\fancypagestyle{empty}{}
\warningExpl{\textbf{This document is a work in progress.}}

\begin{abstract}
This document describes how to use Biber data annotations in \BibLaTeX{} entries, how to use this annotation for \gls{CReDiT} data, and gives some recommendations for generating fake \BibLaTeX{} entries.

\warningExpl{{\mystixmathfont ☡} This is a very deep dive that is probably only of interest to a \textit{very} small number of readers.}
\end{abstract}
\cleardoublepage

%%%%%%%%%%%%%%%%%%%%%%%%%%%%%%%%%%%%%%%%%%%%%%%%%%%%%%%%%%%%%%%%%%%%%%
% Enable the following command the first time you compile this in a new session to give all of the time quanta to set up the font cache.
%% If the font cache is not built, you are likely to get a timeout.
%\end{document}

% Include Table of Contents (TOC) ---
\fancypagestyle{plain}{}
\renewcommand{\sectionmark}[1]{ \markboth{#1}{}}
\tableofcontents
\markboth{\contentsname}{}
\cleardoublepage
 
\setcounter{page}{1}
\pagenumbering{arabic}
\glsresetall
This document describes how to use Biber data annotations in \BibLaTeX{} entries (\Cref{sec:furtherBiberAnnotation}), how to use this annotation for \gls{CReDiT} data (\Cref{sec:annotationsForCReDiT}), and gives some recommendations for generating fake \BibLaTeX{} entries (\Cref{sec:generatingFakeBibLaTeXEntries}).

\section{Roadmap: How to Read This Guide}

Depending on your immediate goals with the \gls{KTH} thesis template, you can navigate this document using the following tracks:
\begin{enumerate}[leftmargin=*, label=\textbf{Track \arabic*}, ref=Track \arabic*]
    \item \textbf{Quick Implementation (Just the ORCID \& Icon Snippets)}
    \begin{itemize}
    \item Go directly to \Cref{sec:furtherBiberAnnotation}
    \item 
Review \Cref{lst:biberAuthorAnnotation} for the `\enquote{author+an} syntax, and jump to \Cref{sec:orcidlinksMethodA} or \Cref{sec:orcidlinksMethodB} for the ready-to-use Lua and LaTeX hooks that automatically render clickable green ORCID badges.
    \end{itemize}

    \item \textbf{Automated Workflows \& the CReDiT Python Pipeline}
    \begin{itemize}
    \item  Focus on \Cref{sec:annotationsForCReDiT}
    
    \item  Review Pattern A (\Cref{lst:biberCreditPatternA}) for granular, degree-based contribution mapping.
       
    \item Skip to \Cref{sec:integratingAnnotationsWithCurrentWorkflows} to understand how external tools like \enquote{\fname{CReDiT\_Matrix\_Wizard.py}} and \enquote{\fname{publications\_map.json}} synchronize with GitHub Actions to build contribution matrices automatically.
    \end{itemize}
    
\item \textbf{Advanced Architecture \& Metadata Overlays}
    \begin{itemize} 
    \item Read \Cref{sec:usingAnnotationsWithReferenceManagers} to understand why standard reference managers (such as Zotero or JabRef) strip data annotations and how the \textbf{Split-File Overlay} \enquote{\fname{references.bib}} + \enquote{\fname{references-annotations.bib}} solves this.
    \item Explore \Cref{sec:otherUsesForBiberDataAnotations} for advanced applications such as Open Access compliance tracking, math-in-title normalization, and handling mega-collaborations.
    \end{itemize}
\end{enumerate}


\section[Using Biber data annotations in BibLaTeX entries]{Using Biber data annotations in BibLaTeX\linebreak[4] entries}
\label{sec:furtherBiberAnnotation}

Biber's data annotation syntax enables bibliography entries to carry rich, structured metadata attached directly to individual author or field items without modifying the underlying \BibLaTeX{} data model. \Cref{lst:biberAuthorAnnotation} shows an example with three authors, each of whom has an \gls{ORCID} identifier, \ie an ORCID iD. Additionally, the first two authors are noted as having contributed equally (\eg to designate joint first authorship).

As described in §3.7 of the \texttt{biblatex} package manual, data annotations enable an author to attach semantic information to information in a bibliography. The field name suffix \enquote{+an} is a user-definable marker that this field contains an annotation. After this marker is an optional named marker (by default \enquote{default}). For example, in \Cref{lst:biberAuthorAnnotation} one can indicate the Swedish title for a work by adding: \texttt{title+an:swedish} and the Swedish title.

\Needspace*{18\baselineskip}
\begin{lstlisting}[language={BibTeX}, breaklines=true, basicstyle=\small\ttfamily, caption={{Biber author annotations for ORCID and equal contribution}}, label=lst:biberAuthorAnnotation]
@article{doe2025collaborative,
  author    = {Maguire, Jr., Gerald Q. and Smith, Jane and Green, Alex},
  author+an = {
    1:orcid="0000-0002-6066-7469";
    1:equal="true";
    2:orcid="0000-0001-2345-6789";
    2:equal="true";
    3:orcid="0000-0003-9876-5432";
  },
  title     = {Collaborative Systems in Higher Education},
  title+an:swedish = {Samarbetssystem inom högre utbildning},
  journal   = {Journal of Academic Stewardship},
  year      = {2025},
  volume    = {3},
  number    = {2},
  pages     = {112--128},
  doi       = {10.1000/182},
  issn       = {0000-0000},
  placeholder={Placeholder fake reference for use in examples}
}
\end{lstlisting}

Note that each line of the author annotation(shown in \Cref{lst:biberAuthorAnnotation}) has a prefix of the form `<number>:' followed by a name, an equal sign (`='), and a quoted value. Internally, this information is encoded into \textit{parts} and stored by \texttt{biber} in the \texttt{`.bbl'} file. Note also the use of the \texttt{placeholder} field to explicitly identify this as a fake reference to be used simply for the purposes of examples.

\subsection{BibLaTeX Part Annotation Signature}
Programmatic access to the annotation is done via:
\begin{small}
\begin{verbatim}
\texttt{\textbackslash haspartannotation[<field>][<name>][<item>]\{<part>\}\{\dots\}\{\dots\}}

\texttt{\textbackslash getpartannotation[<field>][<name>][<item>]\{<part>\}}
\end{verbatim}
\end{small}

The three optional brackets define the bibliographic coordinates (`[author][default][<idx>]`), while the part slug is passed as the mandatory brace argument (`\{<part>\}`). In the example above, the slug will be \enquote{orcid} or \enquote{equal}.

To render these annotations in print, one can define \BibLaTeX{} formatting hooks (as shown in \Cref{lst:biberHookExample1A} on page \pageref{lst:biberHookExample1A}). You can test for active annotations using \texttt{\textbackslash haspartannotation} and retrieve values with \texttt{\textbackslash getpartannotation}. In this code, equal authorship is indicated with a superscript asterisk, while the name gets an  \texttt{\textbackslash orcidlink} hyperlink, \ie beside the name is the official green, circular \gls{ORCID} icon as an actionable, live hyperlink targeting the author's \gls{ORCID} record page. The icon itself is drawn as a resolution-independent vector path using TikZ/PDF literals rather than a raster image.

\Cref{lst:biberHookIconOrText} shows how you can generate the green, circular \gls{ORCID} icon (requires use of the \texttt{orcidlink} package) or a tiny bold font green string saying \enquote{iD}. Note that \texttt{\textbackslash orcidlink\{\}} provides its own spacing after the name and before the icon.

\begin{lstlisting}[style=latexExampleForAuthors, escapechar=|, basicstyle=\footnotesize, caption={{ORCID icon or colored text}}, label=lst:biberHookIconOrText]
\usepackage{orcidlink}

\newcommand*{\renderorcidbadge}[1]{%
\orcidlink{#1}%
}

\newcommand*{\renderorcidbadgeText}[1]{%
\,\href{https://orcid.org/#1}{\textcolor[HTML]{A6CE39}{\tiny\textbf{iD}}}%
}
\end{lstlisting}

\Cref{sec:orcidlinksMethodA} and \Cref{sec:orcidlinksMethodB} demonstrate two different methods of getting an author's ORCID iD and displaying the badge with a hyperlink to the author's \gls{ORCID} landing page. First, we shall note some common features of each method:
both create a \TeX{} box (called \texttt{\textbackslash orcidbox}) to render text into, both feature a new command to render the \gls{ORCID} badge, each defines their own \texttt{render\_orcid\_from\_box()} Lua function, and within a group \first each defines a local \texttt{\textbackslash mkbibnamefamily}\footnote{One can also renew the command \texttt{\textbackslash mkbibnamegiven} for markup of the author's first name. Similarly, the other parts of a name can be utilized: \texttt{\textbackslash mkbibnameprefix} and
\texttt{\textbackslash mkbibnamesuffix}.} that is used when rendering a family name in a bibliography, \Second sets the values of the \texttt{maxnames} and \texttt{minnames} counters to be used for the next bibliographic citation (via 
\texttt{\textbackslash AtNextCite}), and \third uses \texttt{\textbackslash fullcite\{\}} to render a given bibliographic entry.

The two methods differ in their definition of the \texttt{render\_orcid\_from\_box()} or \texttt{render\_orcid\_clean()} Lua function and their own \texttt{\textbackslash renderorcidbadgeX} function. The main difference is how they walk the contents, \ie the rendered nodes in the box. Some of the questions that might come to mind are: Why do we need to render to a box? and Why do we need to walk the contents of this box? The answer to both questions is that the \texttt{\textbackslash getpartannotation} does not return a character string, but \texttt{\textbackslash href} or \texttt{\textbackslash orcidlink} operate in URL-string mode; hence, it requires an evaluated, literal string. By setting the result of \texttt{\textbackslash getpartannotation} in a box, we can now examine its contents and turn it into a string. The code also takes advantage of the fixed format of an ORCID iD to avoid problems with hyphenation and discretionary hyphenation (this is a deep morass in the \LaTeX{} engine used, the font used, many different types of hyphens, and even ligature processing). The two methods utilize different methods of walking the nodes of the box's contents:
\begin{itemize}
    \item Method A (Deep Recursive Walk): Walks into the sub-lists of discretionary nodes (\texttt{n.replace} and \texttt{n.pre}) directly, attempting to extract the literal glyphs embedded inside hyphenation alternatives.

    \item Method B (Box Descent with Direct Hyphen Substitution): Explicitly descends into nested horizontal and vertical box lists (HLIST and VLIST), treats any encountered discretionary node directly as a hyphen (\texttt{table.insert(chars, \enquote{-})}), and relies on the subsequent 16-character normalization pass to format the canonical 4-4-4-4 structure regardless of font shaping effects.
\end{itemize}


\subsection{Method A for inserting ORCID links}
\label{sec:orcidlinksMethodA}
As an example of this for \textbackslash fullcite\{doe2025collaborative\} using the code in \Cref{lst:biberHookExample1A} is:
\newbox\orcidbox

\newcommand*{\renderorcidbadge}[1]{%
\orcidlink{#1}%
}

\begin{luacode*}
function render_orcid_from_box()
  local box_head = tex.box.orcidbox.head
  if not box_head then return end

  local chars = {}
  local GLYPH = node.id("glyph")
  local DISC  = node.id("discretionary")

  local function walk(head)
    for n in node.traverse(head) do
      if n.id == GLYPH then
        table.insert(chars, utf8.char(n.char))
      elseif n.id == DISC then
        if n.replace then
          walk(n.replace)
        elseif n.pre then
          walk(n.pre)
        else
          table.insert(chars, "-")
        end
      elseif n.head then
        walk(n.head)
      end
    end
  end

  walk(box_head)
  local raw_orcid = table.concat(chars)

  -- Filter out anything except digits and X/x (hyphen placed at the very end of class)
  local base = ""
  for c in raw_orcid:gmatch("[0-9xX]") do
    base = base .. c
  end

  local formatted_orcid = ""
  if #base == 16 then
    -- Normalize the 16th character (check digit) to uppercase X if applicable
    local prefix = base:sub(1, 15)
    local checksum = base:sub(16, 16):upper()
    local full = prefix .. checksum

    formatted_orcid = string.format("%s-%s-%s-%s",
      full:sub(1, 4),
      full:sub(5, 8),
      full:sub(9, 12),
      full:sub(13, 16)
    )
  else
    formatted_orcid = raw_orcid
  end

  if formatted_orcid and formatted_orcid ~= "" then
    tex.sprint(token.create("renderorcidbadge"), "{", formatted_orcid, "}")
  end
end
\end{luacode*}

\begingroup
\renewcommand*{\mkbibnamefamily}[1]{%
  #1%
  \edef\curridx{\the\value{listcount}}%
  \haspartannotation[author][default][\curridx]{equal}{%
    \textsuperscript{*}%
  }{}%
  \haspartannotation[author][default][\curridx]{orcid}{%
    \global\setbox\orcidbox=\hbox{\getpartannotation[author][default][\curridx]{orcid}}%
    \directlua{render_orcid_from_box()}%
  }{}%
}

\AtNextCite{%
  \defcounter{maxnames}{99}%
  \defcounter{minnames}{99}%
}

\fullcite{doe2025collaborative}
\endgroup
\medskip

Note that this reference has a red warning \enquote{[Unresolved Placeholder: Placeholder fake reference for use in examples]} to explicitly indicate that it is \emph{not} a real reference. This markup is done via the field declaration in \fname{ext-eprint-dbx.tex} and realized by the check \texttt{\textbackslash iffieldundef\{placeholder\}} invoked by \texttt{\textbackslash AtEveryCitekey} (executes when a citation key is processed in text, \eg \texttt{\textbackslash cite}, \texttt{\textbackslash parencite}, \texttt{\textbackslash textcite}, and \texttt{\textbackslash fullcite}) and \texttt{\textbackslash AtEveryBibitem} (executes inside the bibliography list, \ie \texttt{\textbackslash printbibliography}). Because \texttt{\textbackslash fullcite} internally expands using bibliography drivers, there is no double-printing or omission of the warning on standalone citations in the log file.

\Needspace*{10\baselineskip} % to prevent just a tiny portion being shown at the bottom of a page
\begin{lstlisting}[style=latexExampleForAuthors, escapechar=|, basicstyle=\footnotesize, caption={{Example of a \LaTeX{} hook for equal contribution and ORCID icons}}, label=lst:biberHookExample1A]
\newbox\orcidbox

\newcommand*{\renderorcidbadge}[1]{%
\orcidlink{#1}%
}

\begin{luacode*}
function render_orcid_from_box()
  local box_head = tex.box.orcidbox.head
  if not box_head then return end

  local chars = {}
  local GLYPH = node.id("glyph")
  local DISC  = node.id("discretionary")

  local function walk(head)
    for n in node.traverse(head) do
      if n.id == GLYPH then
        table.insert(chars, utf8.char(n.char))
      elseif n.id == DISC then
        if n.replace then
          walk(n.replace)
        elseif n.pre then
          walk(n.pre)
        else
          table.insert(chars, "-")
        end
      elseif n.head then
        walk(n.head)
      end
    end
  end

  walk(box_head)
  local raw_orcid = table.concat(chars)

  -- Filter out anything except digits and X/x (hyphen placed at the very end of class)
  local base = ""
  for c in raw_orcid:gmatch("[0-9xX]") do
    base = base .. c
  end

  local formatted_orcid = ""
  if #base == 16 then
    -- Normalize the 16th character (check digit) to uppercase X if applicable
    local prefix = base:sub(1, 15)
    local checksum = base:sub(16, 16):upper()
    local full = prefix .. checksum

    formatted_orcid = string.format("%s-%s-%s-%s",
      full:sub(1, 4),
      full:sub(5, 8),
      full:sub(9, 12),
      full:sub(13, 16)
    )
  else
    formatted_orcid = raw_orcid
  end

  if formatted_orcid and formatted_orcid ~= "" then
    tex.sprint(token.create("renderorcidbadge"), "{", formatted_orcid, "}")
  end
end
\end{luacode*}

\begingroup
\renewcommand*{\mkbibnamefamily}[1]{%
  #1%
  \edef\curridx{\the\value{listcount}}%
  \haspartannotation[author][default][\curridx]{equal}{%
    \textsuperscript{*}%
  }{}%
  \haspartannotation[author][default][\curridx]{orcid}{%
    \global\setbox\orcidbox=\hbox{\getpartannotation[author][default][\curridx]{orcid}}%
    \directlua{render_orcid_from_box()}%
  }{}%
}

\AtNextCite{%
  \defcounter{maxnames}{99}%
  \defcounter{minnames}{99}%
}

\fullcite{doe2025collaborative}
\endgroup
\end{lstlisting}




\subsection{Method B for inserting ORCID links}
\label{sec:orcidlinksMethodB}
Another example of inserting ORCID links for \textbackslash fullcite\{doe2025collaborative\} using the code in \Cref{lst:biberHookExample1B} is:
\newbox\orcidbox

\newcommand*{\renderorcidbadgeB}[1]{%
\orcidlink{#1}%
}

\begin{luacode*}
function render_orcid_clean()
  local box_head = tex.box.orcidbox.head
  if not box_head then return end

  local chars = {}
  local GLYPH = node.id("glyph")
  local DISC  = node.id("discretionary")
  local HLIST = node.id("hlist")
  local VLIST = node.id("vlist")

  local function walk(head)
    for n in node.traverse(head) do
      if n.id == GLYPH then
        table.insert(chars, utf8.char(n.char))
      elseif n.id == DISC then
        table.insert(chars, "-")
      elseif n.id == HLIST or n.id == VLIST then
        if n.head then walk(n.head) end
      end
    end
  end

  walk(box_head)
  local raw_orcid = table.concat(chars)

  -- 1. Extract only alphanumeric characters (digits and X/x)
  local alphanum = ""
  for c in raw_orcid:gmatch("[0-9xX]") do
    alphanum = alphanum .. c
  end

  -- 2. Format into standard ORCID 4-4-4-4 hyphenated layout
  local final_orcid = ""
  if #alphanum == 16 then
    local prefix = alphanum:sub(1, 15)
    local checksum = alphanum:sub(16, 16):upper()
    local full = prefix .. checksum

    final_orcid = string.format("%s-%s-%s-%s",
      full:sub(1, 4),
      full:sub(5, 8),
      full:sub(9, 12),
      full:sub(13, 16)
    )
  else
    final_orcid = raw_orcid
  end

  if final_orcid and final_orcid ~= "" then
    tex.sprint(token.create("renderorcidbadgeB"), "{", final_orcid, "}")
  end
end
\end{luacode*}

\begingroup
\renewcommand*{\mkbibnamefamily}[1]{%
  #1%
  \edef\curridx{\the\value{listcount}}%
  \haspartannotation[author][default][\curridx]{equal}{%
    \textsuperscript{*}%
  }{}%
  \haspartannotation[author][default][\curridx]{orcid}{%
    \global\setbox\orcidbox=\hbox{\getpartannotation[author][default][\curridx]{orcid}}%
    \directlua{render_orcid_clean()}%
  }{}%
}

\AtNextCite{%
  \defcounter{maxnames}{99}%
  \defcounter{minnames}{99}%
}

\fullcite{doe2025collaborative}
\endgroup
\begin{lstlisting}[style=latexExampleForAuthors, escapechar=|, basicstyle=\footnotesize, caption={{Example of a \LaTeX{} hook for equal contribution and ORCID icons}}, label=lst:biberHookExample1B]
\usepackage{biblatex}
\usepackage{orcidlink}

\newbox\orcidbox

\newcommand*{\renderorcidbadgeB}[1]{%
\orcidlink{#1}%
}

\begin{luacode*}
function render_orcid_clean()
  local box_head = tex.box.orcidbox.head
  if not box_head then return end

  local chars = {}
  local GLYPH = node.id("glyph")
  local DISC  = node.id("discretionary")
  local HLIST = node.id("hlist")
  local VLIST = node.id("vlist")

  local function walk(head)
    for n in node.traverse(head) do
      if n.id == GLYPH then
        table.insert(chars, utf8.char(n.char))
      elseif n.id == DISC then
        table.insert(chars, "-")
      elseif n.id == HLIST or n.id == VLIST then
        if n.head then walk(n.head) end
      end
    end
  end

  walk(box_head)
  local raw_orcid = table.concat(chars)

  -- 1. Extract only alphanumeric characters (digits and X/x)
  local alphanum = ""
  for c in raw_orcid:gmatch("[0-9xX]") do
    alphanum = alphanum .. c
  end

  -- 2. Format into standard ORCID 4-4-4-4 hyphenated layout
  local final_orcid = ""
  if #alphanum == 16 then
    local prefix = alphanum:sub(1, 15)
    local checksum = alphanum:sub(16, 16):upper()
    local full = prefix .. checksum

    final_orcid = string.format("%s-%s-%s-%s",
      full:sub(1, 4),
      full:sub(5, 8),
      full:sub(9, 12),
      full:sub(13, 16)
    )
  else
    final_orcid = raw_orcid
  end

  if final_orcid and final_orcid ~= "" then
    tex.sprint(token.create("renderorcidbadgeB"), "{", final_orcid, "}")
  end
end
\end{luacode*}

\begingroup
\renewcommand*{\mkbibnamefamily}[1]{%
  #1%
  \edef\curridx{\the\value{listcount}}%
  \haspartannotation[author][default][\curridx]{equal}{%
    \textsuperscript{*}%
  }{}%
  \haspartannotation[author][default][\curridx]{orcid}{%
    \global\setbox\orcidbox=\hbox{\getpartannotation[author][default][\curridx]{orcid}}%
    \directlua{render_orcid_clean()}%
  }{}%
}

\AtNextCite{%
  \defcounter{maxnames}{99}%
  \defcounter{minnames}{99}%
}

\fullcite{doe2025collaborative}
\endgroup
\end{lstlisting}

\subsection{Using annotation to store Contributor Roles Taxonomy (CReDiT) data}
\label{sec:annotationsForCReDiT}

Currently, there are no hardcoded built-in keyword annotations reserved specifically for \gls{CReDiT} (Contributor Roles Taxonomy, ANSI/NISO Z39.104-2022~\cite{niso_z39_104_2022}). However, standard canonical identifiers using the \enquote{Simplified Name} form can be utilized directly as annotation keys, as outlined in \Cref{tab:CReDiT_annotation}. The degree values of these can be: \enquote{lead}, \enquote{equal} (=true), or \enquote{supporting}:
\begin{description}[labelwidth =\widthof{\textbf{supporting}}, leftmargin = !]
    \item[lead] The contributor took primary responsibility for the facet of the work.

    \item[equal] The contributor shared substantial, balanced responsibility alongside one or more co-authors.

    \item[supporting] The contributor provided supplementary contribution, operational assistance, or targeted input.
\end{description}
Absence of the key implies no attributed contribution.

\begin{table}[!htp]
  \centering
  \caption{CReDiT roles for author data annotations in \BibLaTeX with their Simplified Name, \ie annotation key }
  \label{tab:CReDiT_annotation}
  \begin{tabular}{L{5.5cm} L{8.5cm}}
    \hline
    \textbf{CReDiT Role} & \textbf{Simplified Name} \\
    \hline
    Conceptualization          & \texttt{conceptualization} \\
    Data curation              & \texttt{data-curation} \\
    Formal analysis            & \texttt{formal-analysis}\\
    Funding acquisition        & \texttt{funding-acquisition} \\
    Investigation              & \texttt{investigation}\\
    Methodology                & \texttt{methodology} \\
    Project administration     & \texttt{project-administration} \\
    Resources                  & \texttt{resources} \\
    Software                   & \texttt{software} \\
    Supervision                & \texttt{supervision} \\
    Validation                 & \texttt{validation} \\
    Visualization              & \texttt{visualization} \\
    Writing -- original draft  & \texttt{writing-original-draft} \\
    Writing -- review \& editing & \texttt{writing-review-editing}\\
    \hline
  \end{tabular}
\end{table}
\FloatBarrier

There are two primary paradigms for encoding \gls{CReDiT} roles using \texttt{biber} annotations, as shown in \Cref{lst:biberCreditPatternA} and \Cref{lst:biberCreditPatternB}.

\begin{lstlisting}[language={BibTeX}, breaklines=true, basicstyle=\small\ttfamily, caption={{Pattern A: Role-as-Key (Fine-grained with degrees of contribution)}}, label=lst:biberCreditPatternA]
@article{maguire2026pipeline,
  author    = {Maguire Jr., Gerald Q. and Smith, Jane},
  author+an = {
    1:orcid="0000-0002-6066-746X";
    1:conceptualization="lead";
    1:software="lead";
    1:writing-original-draft="lead";
    2:orcid="0000-0001-2345-6789";
    2:supervision="lead";
    2:writing-review-editing="supporting"
  },
  title     = {Automated Metadata Integrity in Academic Publishing},
  journal   = {Journal of Academic Stewardship},
  year      = {2026},
  volume    = {4},
  number    = {1},
  pages     = {45--60},
  doi       = {10.1000/182},
  placeholder={Placeholder fake reference for use in examples}
}
\end{lstlisting}

\begin{lstlisting}[language={BibTeX}, breaklines=true, basicstyle=\small\ttfamily, caption={{Pattern B: Single credit container (Comma-separated token string)}}, label=lst:biberCreditPatternB]
@article{maguire2026pipeline,
  author    = {Maguire, Jr., Gerald Q. and Smith, Jane},
  author+an = {
    1:orcid="0000-0002-6066-746X"; 1:credit="Conceptualization, Software, Writing - original draft";
    2:orcid="0000-0001-2345-6789"; 2:credit="Supervision, Writing - review & editing"
  },
  title     = {Automated Metadata Integrity in Academic Publishing},
  journal   = {Journal of Academic Stewardship},
  year      = {2026},
  placeholder = {A fake example for documentation}
}
\end{lstlisting}

Pattern A is the best method for \gls{KTH} doctoral thesis students because it makes it easy to state the values associated with a role.

\Needspace*{22\baselineskip}
\begin{lstlisting}[style=latexExampleForAuthors, escapechar=|, basicstyle=\footnotesize, caption={{Example \LaTeX{} hook for extracting and rendering contributor roles}}, label=lst:biberHookExample2]
\usepackage{biblatex}

% Macro to iterate through author list and render contributor role statements
\newbibmacro*{author:credit:roles}{%
  \begingroup
  \def\do{%
    \haspartannotation[author][default][\currentname]{credit}%
      {\par\small\textbf{\namepartgiven} \textbf{\namepartfamily}: \getpartannotation[author][default][\currentname]{credit}}%
      {}%
  }%
  \printnames[creditnames]{author}%
  \endgroup
}

% Append contribution details to extended bibliography listings
\renewbibmacro*{finentry}{%
  \usebibmacro{author:credit:roles}%
  \finentry%
}
\end{lstlisting}

This setup allows automated Python/\gls{CI} validation pipelines to extract and cross-check contributor statements directly from \fname{.bib} files without manual re-entry.

Because \texttt{\textbackslash renewbibmacro*} is dynamically scoped in \LaTeX, it can be redefined specifically for the author contribution section and either restored afterward or restricted within a local group.

As shown in \Cref{lst:biberHookExampleLocalGroupScoping}, wrapping the contribution section inside a \TeX\ group (\texttt{\textbackslash begingroup \ldots \textbackslash endgroup}) ensures the modification automatically discards itself when the scope closes, leaving the main thesis bibliography unaffected:

\Needspace*{18\baselineskip}
\begin{lstlisting}[style=latexExampleForAuthors, escapechar=|, basicstyle=\footnotesize, caption={{Local group scoping for CReDiT bibliography output}}, label=lst:biberHookExampleLocalGroupScoping]
\chapter*{Author Contributions}

\begingroup
% 1. Redefine finentry locally to append CReDiT statements
\renewbibmacro*{finentry}{%
  \usebibmacro{author:credit:roles}%
  \finentry%
}

\textbf{Paper I}
\fullcite{maguire2026pipeline}

\textbf{Paper II}
\fullcite{doe2025collaborative}
\endgroup

% Outside the group, finentry reverts to its original standard definition
\end{lstlisting}

\Needspace*{20\baselineskip}
The following generates some example output when using Approach A for the \gls{CReDiT} information in the \fname{.bib} file:
\begingroup
% 1. Helper mapping from simplified kebab-case slug to display text
\newcommand{\formatcreditrole}[1]{%
  \ifcsdef{credit@#1}{\csuse{credit@#1}}{\detokenize{#1}}%
}
\csdef{credit@conceptualization}{Conceptualization}
\csdef{credit@data-curation}{Data curation}
\csdef{credit@formal-analysis}{Formal analysis}
\csdef{credit@funding-acquisition}{Funding acquisition}
\csdef{credit@investigation}{Investigation}
\csdef{credit@methodology}{Methodology}
\csdef{credit@project-administration}{Project administration}
\csdef{credit@resources}{Resources}
\csdef{credit@software}{Software}
\csdef{credit@supervision}{Supervision}
\csdef{credit@validation}{Validation}
\csdef{credit@visualization}{Visualization}
\csdef{credit@writing-original-draft}{Writing -- original draft}
\csdef{credit@writing-review-editing}{Writing -- review \& editing}
\csdef{credit@equal}{Equal contribution}

% 2. List of CReDiT simplified slugs to check
\newcommand*{\CreditSlugList}{%
  conceptualization,data-curation,formal-analysis,funding-acquisition,%
  investigation,methodology,project-administration,resources,software,%
  supervision,validation,visualization,writing-original-draft,writing-review-editing,equal%
}

% 3. Check and print role:
%    Syntax: \haspartannotation[<field>][<annotation_name>][<item_index>]{<part_name>}{<true>}{<false>}
%    Syntax: \getpartannotation[<field>][<annotation_name>][<item_index>]{<part_name>}
\newcommand*{\checkandprintrole}[2]{%
  % #1 = author integer, #2 = slug
  \haspartannotation[author][default][#1]{#2}{%
    \par\hspace*{1.5em}\formatcreditrole{#2}: \textit{\getpartannotation[author][default][#1]{#2}}%
  }{}%
}

% 4. Name format driver
\DeclareNameFormat{creditnames}{%
  \edef\currentauthornum{\the\value{listcount}}%
  \par\smallskip\noindent\textbf{\namepartgiven} \textbf{\namepartfamily}%
  \def\do##1{\checkandprintrole{\currentauthornum}{##1}}%
  \expandafter\docsvlist\expandafter{\CreditSlugList}%
}

% 5. Bibmacro: forces evaluation of all authors
\newbibmacro*{author:credit:roles}{%
  \begingroup
  \defcounter{maxnames}{99}%
  \defcounter{minnames}{99}%
  \par\medskip\noindent\textbf{Author Contributions:}%
  \printnames[creditnames]{author}%
  \endgroup
}

% 6. Hook into finentry
\renewbibmacro*{finentry}{%
  \finentry
  \usebibmacro{author:credit:roles}%
}

\textbf{Paper I}\\
\fullcite{maguire2026pipeline}

\Needspace*{20\baselineskip}
\textbf{Paper II}\\
\fullcite{doe2025collaborative}
\endgroup



\subsection{Architectural Takeaways for Your Pipeline}
\begin{itemize}
    \item \textbf{biber \enquote{.bib} Annotation Delimiters:} Item annotations must be delimited by semicolons (`;`), not commas (`,`), with the item index explicit on each part (\eg `1:conceptualization="lead"; 1:software="lead"`).
    \item \textbf{Kebab-Case Slugs:} Adopting the standard \gls{CReDiT} simplified names (\eg `writing-original-draft`, `data-curation`) avoids unescaped underscore collisions in TeX math mode and simplifies downstream \gls{JSON} serialization.
    \item \textbf{BibLaTeX Part Annotation Signature:}
\begin{small}
\begin{verbatim}
\texttt{\textbackslash haspartannotation[<field>][<name>][<item>]\{<part>\}\{\dots\}\{\dots\}}


\texttt{\textbackslash getpartannotation[<field>][<name>][<item>]\{<part>\}}
\end{verbatim}
\end{small}

The three optional brackets define the bibliographic coordinates (`[author][default][<idx>]`), while the part slug is passed as the mandatory brace argument (`\{<part>\}`).
    \item \textbf{Index Isolation:} Freezing `\textbackslash the\textbackslash value{listcount}` via `\textbackslash edef` before triggering `\textbackslash docsvlist` prevents the inner loop from clobbering the active author index.
\end{itemize} 


If you prefer to redefine macros globally across document sections, redefine \texttt{finentry} before the papers list and reset it before calling \texttt{\textbackslash printbibliography}:

\Needspace*{16\baselineskip}
\begin{lstlisting}[style=latexExampleForAuthors, escapechar=|, basicstyle=\footnotesize, caption={{Sequential document-level redefinition of finentry}}, label=lst:biberHookExampleSequentialRedefinition]
% In the Author Contributions front matter:
\renewbibmacro*{finentry}{%
  \usebibmacro{author:credit:roles}%
  \finentry%
}

\fullcite{maguire2026pipeline}

% -------------------------------------------------------------
% Later, in the back matter before the main thesis bibliography:
% -------------------------------------------------------------

\renewbibmacro*{finentry}{\finentry}

\printbibliography[heading=bibintoc, title={References}]
\end{lstlisting}

For granular control when citing individual papers with contributor details without modifying global hooks, encapsulate the scoped redefinition inside a custom macro:

\Needspace*{12\baselineskip}
\begin{lstlisting}[style=latexExampleForAuthors, escapechar=|, basicstyle=\footnotesize, caption={{Defining a dedicated \textbackslash fullcitecredit command}}, label=lst:biberfullcitecreditCommand]
% Dedicated command encapsulating scoped CReDiT redefinition
\newcommand{\fullcitecredit}[1]{%
  \begingroup
  \renewbibmacro*{finentry}{%
    \usebibmacro{author:credit:roles}%
    \finentry%
  }%
  \fullcite{#1}%
  \endgroup
}
\end{lstlisting}

\BibLaTeX{} provides the author index context automatically via \texttt{\textbackslash currentname} (which expands to the current name list position, such as 1, 2, \etc.) when stepping through an author list with \texttt{\textbackslash indexnames} or custom name formatters.

Define the slug-to-text mappings with \texttt{\textbackslash csdef}, as shown in \Cref{lst:biberCreditRolesIterate}. Using \texttt{etoolbox} list iteration (\texttt{\textbackslash docsvlist}), you can step through each author and test each of the 14 standard \gls{CReDiT} slugs systematically. Note that the code in \Cref{lst:biberCreditRolesIterate} should be placed in the preamble, except for \enquote{6. Hook into finentry} which only needs to be invoked (within a \texttt{\textbackslash begingroup} \dots \texttt{\textbackslash endgroup}) when you need it for listing the author contributions in your thesis.

\Needspace*{21\baselineskip}
\begin{lstlisting}[style=latexExampleForAuthors, escapechar=|, basicstyle=\footnotesize, caption={{Iterate through slugs}}, label=lst:biberCreditRolesIterate]
\usepackage{biblatex}
\usepackage{etoolbox}

% 1. Helper mapping from simplified kebab-case slug to display text
\newcommand{\formatcreditrole}[1]{%
  \ifcsdef{credit@#1}{\csuse{credit@#1}}{\detokenize{#1}}%
}
\csdef{credit@conceptualization}{Conceptualization}
\csdef{credit@data-curation}{Data curation}
\csdef{credit@formal-analysis}{Formal analysis}
\csdef{credit@funding-acquisition}{Funding acquisition}
\csdef{credit@investigation}{Investigation}
\csdef{credit@methodology}{Methodology}
\csdef{credit@project-administration}{Project administration}
\csdef{credit@resources}{Resources}
\csdef{credit@software}{Software}
\csdef{credit@supervision}{Supervision}
\csdef{credit@validation}{Validation}
\csdef{credit@visualization}{Visualization}
\csdef{credit@writing-original-draft}{Writing -- original draft}
\csdef{credit@writing-review-editing}{Writing -- review \& editing}
\csdef{credit@equal}{Equal contribution}

% 2. List of CReDiT simplified slugs to check
\newcommand*{\CreditSlugList}{%
  conceptualization,data-curation,formal-analysis,funding-acquisition,%
  investigation,methodology,project-administration,resources,software,%
  supervision,validation,visualization,writing-original-draft,writing-review-editing,equal%
}

% 3. Check and print role:
\newcommand*{\checkandprintrole}[2]{%
  \haspartannotation[author][default][#1]{#2}{%
    \par\hspace*{1.5em}\formatcreditrole{#2}: \textit{\getpartannotation[author][default][#1]{#2}}%
  }{}%
}

% 4. Name format driver
\DeclareNameFormat{creditnames}{%
  \edef\currentauthornum{\the\value{listcount}}%
  \par\smallskip\noindent\textbf{\namepartgiven} \textbf{\namepartfamily}%
  \def\do##1{\checkandprintrole{\currentauthornum}{##1}}%
  \expandafter\docsvlist\expandafter{\CreditSlugList}%
}

% 5. Bibmacro: forces evaluation of all authors
\newbibmacro*{author:credit:roles}{%
  \begingroup
  \defcounter{maxnames}{99}%
  \defcounter{minnames}{99}%
  \par\medskip\noindent\textbf{Author Contributions:}%
  \printnames[creditnames]{author}%
  \endgroup
}

% 6. Hook into finentry
\renewbibmacro*{finentry}{%
  \finentry
  \usebibmacro{author:credit:roles}%
}
\end{lstlisting}

In \BibLaTeX, whenever a citation or bibliography driver is executing (such as inside \texttt{\textbackslash AtEveryBibitem}, \texttt{\textbackslash AtEveryCitekey}, \texttt{\textbackslash DeclareBibliographyDriver}, \texttt{\textbackslash fullcite}, or inside a bibmacro like \texttt{finentry}), the active entry key is automatically set as the current context.

When writing custom reports, contribution matrices, or diagnostic validation routines outside the standard citation loop, \BibLaTeX{} provides several mechanisms to inspect and manipulate the active citation context:
\begin{enumerate}
    \item Standard Built-in Key Registers
    \begin{itemize}
        \item \texttt{\textbackslash thefield\{entrykey\}}: Prints or retrieves the active \fname{.bib} entry key (\eg maguire2026pipeline).

        \item \texttt{\textbackslash strfield\{entrykey\}}: Returns the unexpanded string representation of the entry key.

        \item \texttt{\textbackslash currententrykey}: An internal macro holding the literal key of the item currently being processed.
    \end{itemize}
    
    \item Explicitly Setting an Arbitrary Entry as \enquote{Current}
    
If you are outside a standard bibliography loop (for example, in custom TeX/Lua code or arbitrary body text) and want to load a specific entry key into the active \BibLaTeX{} evaluation context without printing a full citation:
\end{enumerate}

\begin{lstlisting}[style=latexExampleForAuthors, escapechar=|, basicstyle=\footnotesize, caption={{Load specific entry key}}, label=lst:biberLoadSpecifcEntryKey]
\begingroup
% Switch the active biblatex context to 'myEntryKey'
\entryclone{myEntryKey}{current_context_key} % Optional clone
% or set the entry data context directly:
\usedatacontext[entrykey=myEntryKey]
...
\endgroup
\end{lstlisting}

More commonly, \BibLaTeX{} provides \texttt{\textbackslash entrydata\{<entrykey>\}\{<code/macros>\}}, which temporarily switches the current entry context to \texttt{<entrykey>} and executes your custom code within that entry's data scope:
\begin{lstlisting}[style=latexExampleForAuthors, escapechar=|, basicstyle=\footnotesize, caption={{Temporary switch to specific entry key}}, label=lst:biberTemporarySwitchSpecifcEntryKey]
\entrydata{maguire2026pipeline}{%
  % Inside this block, 'maguire2026pipeline' is the active entry
  \textbf{Active Key:} \thefield{entrykey}\par
  \textbf{Title:} \printfield{title}\par
  \usebibmacro{author:credit:roles}%
}
\end{lstlisting}

\Needspace*{9\baselineskip}
Example usage to print the current entry key alongside metadata;
\begin{lstlisting}[style=latexExampleForAuthors, escapechar=|, basicstyle=\footnotesize, caption={{Print current entry key with metadata}}, label=lst:biberPrintCurrentEntry]
\newbibmacro*{print:current:entry:info}{%
  \par\noindent
  \textbf{Entry Key:} \thefield{entrykey}\par
  \textbf{Entry Type:} \thefield{entrytype}\par
  \textbf{Active Authors:} \printnames{author}%
}
\end{lstlisting}

\subsection{Listing roles grouped by degree}

Whenever the \texttt{author:credit:roles} macro is invoked during bibliography rendering, \texttt{\textbackslash thefield\{entrykey\}} dynamically evaluates to whatever reference is currently being processed. \Cref{lst:biberCReDiTrolesGroupedByDegree} shows how the roles can be listed grouped by degree. If an author has no roles allocated under a specific degree (such as no Equal roles), that row is omitted entirely.
\Needspace*{22\baselineskip}
\begin{lstlisting}[style=latexExampleForAuthors, escapechar=|, basicstyle=\footnotesize, caption={{Listing CReDiT roles grouped by degree}}, label=lst:biberCReDiTrolesGroupedByDegree]
\usepackage{etoolbox}

% Master list of the 14 standard CReDiT roles
\newcommand*{\CreditRolesList}{%
  conceptualization,%
  data-curation,%
  formal-analysis,%
  funding-acquisition,%
  investigation,%
  methodology,%
  project-administration,%
  resources,%
  software,%
  supervision,%
  validation,%
  visualization,%
  writing-original-draft,%
  writing-review-editing%
}

% Human-readable labels for role slugs
\newcommand*{\formatcreditrole}[1]{%
  \ifcsdef{creditlabel@#1}%
    {\csuse{creditlabel@#1}}%
    {#1}%
}
\csdef{creditlabel@conceptualization}{Conceptualization}
\csdef{creditlabel@data-curation}{Data Curation}
\csdef{creditlabel@formal-analysis}{Formal Analysis}
\csdef{creditlabel@funding-acquisition}{Funding Acquisition}
\csdef{creditlabel@investigation}{Investigation}
\csdef{creditlabel@methodology}{Methodology}
\csdef{creditlabel@project-administration}{Project Administration}
\csdef{creditlabel@resources}{Resources}
\csdef{creditlabel@software}{Software}
\csdef{creditlabel@supervision}{Supervision}
\csdef{creditlabel@validation}{Validation}
\csdef{creditlabel@visualization}{Visualization}
\csdef{creditlabel@writing-original-draft}{Writing -- Original Draft}
\csdef{creditlabel@writing-review-editing}{Writing -- Review \& Editing}

% Formatting template for each degree group
\newcommand*{\printcreditgroup}[2]{%
  \ifcsvoid{#1}{}{%
    \par\hskip 1.5em\relax
    \textit{#2:} \csvtolabel{#1}%
  }%
}

% Comma-separated output helper for role lists
\newcommand*{\csvtolabel}[1]{%
  \begingroup
  \def\do##1{\formatcreditrole{##1}\def\do{, \formatcreditrole{####1}}}%
  \dolistcsloop{#1}%
  \endgroup
}

% Grouped synthesis name-formatting driver
\DeclareNameFormat{creditgrouped}{%
  \begingroup
  % Initialize accumulators for each contribution degree
  \csgdef{credit@lead@\the\value{listcount}}{}%
  \csgdef{credit@equal@\the\value{listcount}}{}%
  \csgdef{credit@supporting@\the\value{listcount}}{}%
  \togglefalse{hascredit}%
  %
  % Iterate through all 14 roles for the current author index
  \renewcommand*{\do}[1]{%
    \haspartannotation[author][##1][\the\value{listcount}]%
      {%
        \toggletrue{hascredit}%
        \edef\curdegree{\getpartannotation[author][##1][\the\value{listcount}]}%
        \ifdefstring{\curdegree}{lead}%
          {\listcsgadd{credit@lead@\the\value{listcount}}{##1}}{}%
        \ifdefstring{\curdegree}{equal}%
          {\listcsgadd{credit@equal@\the\value{listcount}}{##1}}{}%
        \ifdefstring{\curdegree}{supporting}%
          {\listcsgadd{credit@supporting@\the\value{listcount}}{##1}}{}%
      }{}%
  }%
  \docsvlist{\CreditRolesList}%
  %
  % Typeset author and grouped roles only if contributions are registered
  \iftoggle{hascredit}{%
    \par\smallskip\noindent
    \textbf{%
      \ifdefvoid{\namepartgiven}{}{\namepartgiven\ }%
      \ifdefvoid{\namepartfamily}{}{\namepartfamily}%
      \ifdefvoid{\namepartjunior}{}{ \namepartjunior}%
    }%
    \haspartannotation[author][orcid][\the\value{listcount}]%
      {\space\footnotesize(\url{https://orcid.org/\getpartannotation[author][orcid][\the\value{listcount}]})}%
      {}%
    \par\nopagebreak
    \printcreditgroup{credit@lead@\the\value{listcount}}{Lead}%
    \printcreditgroup{credit@equal@\the\value{listcount}}{Equal}%
    \printcreditgroup{credit@supporting@\the\value{listcount}}{Supporting}%
  }{}%
  \endgroup
}

% Replacement master macro
\newbibmacro*{author:credit:roles}{%
  \begingroup
  \newtoggle{hascredit}%
  \printnames[creditgrouped]{author}%
  \endgroup
}
\end{lstlisting}

Using creditlabel@\ldots{} provides strict separation:
\begin{description}[labelwidth =\widthof{\textbf{\texttt{credit@<degree>@<author\_idx>}}}, leftmargin = !]
    \item[\texttt{creditlabel@<role>}] Static translation table mapping raw slug keys to typeset strings.

    \item[\texttt{credit@<degree>@<author\_idx>}] Dynamic runtime data containers accumulating roles during macro evaluation.
\end{description}

\subsection{Using biber data annotations for other purposes}
\label{sec:otherUsesForBiberDataAnotations}

Beyond author attribution and multilingual titles, Biber’s data annotations can attach semantic metadata to citations, authors, and sources without requiring custom \texttt{.dbx} datamodel extensions. Potentially, several applications could be integrated directly into thesis workflows:
\begin{enumerate}
\Needspace*{27\baselineskip}
    \item Thesis Inclusion and Paper Categorization (Compilation Theses)

In Swedish compilation doctoral theses (\foreignlanguage{swedish}{sammanläggningsavhandlingar}), candidates must distinguish between papers formally included in the thesis and secondary publications completed during their studies:

\begin{lstlisting}[language={BibTeX}, breaklines=true, basicstyle=\small\ttfamily,  escapechar=|, caption={{Storing thesis related metadata}}, label=lst:biberThesisMetadata]
@article{maguire2025routing,
  author     = {Maguire, Jr., Gerald Q. and Student, Alice},
  entry+an   = {included="true"; paper="Paper I"},
  title      = {Adaptive Mesh Architectures in 6G Deployments},
  journal    = {IEEE Transactions on Networking},
  year       = {2025},
  placeholder={Placeholder fake reference for use in examples}
}

@inproceedings{student2024early,
  author     = {Student, Alice and Maguire, Jr., Gerald Q.},
  entry+an   = {included="false"; category="other"},
  title      = {Preliminary Bounds on Packet Dispersion},
  booktitle  = {Proceedings of ACM SIGCOMM Posters},
  year       = {2024},
  placeholder={Placeholder fake reference for use in examples}
}
\end{lstlisting}
\begin{description}
    \item[Use in Document] A single \fname{.bib} file can populate both the \textbf{Included Papers} section and the \textbf{Other Publications} appendix by filtering on  \texttt{\textbackslash hasentryannotation\{included\}} inside \texttt{\textbackslash defbibfilter} or \texttt{\textbackslash defbibcheck}.
    \item[Automated Dividers] The value \texttt{paper="Paper I"} can be dynamically fetched to typeset decorative Roman-numeral divider sheets preceding individual reprint \glspl{PDF}.
\end{description}
\Needspace*{20\baselineskip}
    \item Institutional Affiliation and Degree Eligibility Mapping

For co-authored papers, evaluating committees often need to confirm whether an author was affiliated with \gls{KTH} or performed the research under an external industrial partnership:
\begin{lstlisting}[language={BibTeX}, breaklines=true, basicstyle=\small\ttfamily, escapechar=|, caption={{Storing thesis related metadata}}, label=lst:biberThesisAffiliationMetadata]
@article{smith2026telecom,
  author    = {Smith, Jane and Maguire, Jr., Gerald Q. and Larsson, Sven},
  author+an = {
    1:affiliation="kth"; 1:school="EECS";
    2:affiliation="kth"; 2:school="EECS";
    3:affiliation="external"; 3:org="Ericsson Research"
  },
  title     = {Deterministic Latency Guarantees in TSN},
  journal   = {Computer Networks},
  year      = {2026},
  placeholder={Placeholder fake reference for use in examples}
}
\end{lstlisting}
\begin{description}
    \item[Use in Document] Footnotes or superscript markers can identify authors belonging to specific departments or external bodies without cluttering standard citation styles.
    \item[MODS / DiVA Generation] The Python pipeline can inspect \texttt{affiliation} and \texttt{org} to construct \texttt{<affiliation>} and corporate name nodes for \gls{DiVA} without manual disambiguation.
\end{description}

While footnotes or markers to identify authors belonging to specific departments or external bodies might be possible, I think they add too much cognitive load for the reader; therefore, it is better to include these descriptions in the text. Moreover, a better means to input metadata into \gls{DiVA} is needed (potentially, through its new \gls{API}); however, this lies outside the scope of the current work.
\Needspace*{18\baselineskip}
    \item Open Access (OA) \& Archival Compliance Status

Institutional repositories and funding bodies (such as the Swedish Research Council / \foreignlanguage{swedish}{Vetenskapsrådet}) track whether listed papers meet Plan S open-access criteria:
\begin{lstlisting}[language={BibTeX}, breaklines=true, basicstyle=\small\ttfamily,  escapechar=|, caption={{Storing Open Access (OA) \& Archival Compliance Status}}, label=lst:biberOAMetadata]
@article{johnson2026protocols,
  author   = {Johnson, Mark and Maguire, Jr., Gerald Q.},
  entry+an = {oa="gold"; license="cc-by-4.0"; repository="diva:123456"},
  title    = {P4 Acceleration on SmartNIC Clusters},
  journal  = {IEEE Communications Surveys & Tutorials},
  year     = {2026},
  placeholder={Placeholder fake reference for use in examples}
}
\end{lstlisting}

\begin{description}
    \item[Use in Document] Badges (\eg small open-access locks or CC icons) can be rendered beside citations in the list of publications.
    \item[Verification Scripts] A Python/\gls{CI} pre-flight check can flag any included paper annotated as `oa="closed"` or missing a repository identifier.
\end{description}

This looks like a clear area for future work. It can build on the earlier \gls{ORCID} code.
\Needspace*{20\baselineskip}
    \item Reproducibility, Artifact, and Dataset Linking

Conferences and journals increasingly evaluate research artifacts (software, raw datasets, benchmarks). Annotating field-level URLs clarifies the target resource:
\begin{lstlisting}[language={BibTeX}, breaklines=true, basicstyle=\small\ttfamily,  escapechar=|, caption={{Storing Reproducibility, Artifact, and Dataset Linking}}, label=lst:biberRepoducabilityMetadata]
@article{turner2025inference,
  author   = {Turner, Bob and Smith, Jane},
  url      = {https://doi.org/10.1000/182},
  url+an   = {type="publication"},
  eprint   = {10.5281/zenodo.1234567},
  eprint+an = {type="dataset"; badge="zenodo"},
  note+an  = {artifact="acm-reproduced-badge"},
  title    = {cuBLAS Acceleration for Structured Pruning},
  journal  = {ACM Transactions on Computer Systems},
  year     = {2025},
  placeholder={Placeholder fake reference for use in examples}
}
\end{lstlisting}
\begin{description}
    \item[Use in Document] Field annotations on `eprint` or `url` allow custom bibliography macros to display distinct badges (\eg Zenodo dataset icon, GitHub repository badge, or ACM Reproducibility Seals) directly adjacent to the reference.
\end{description}

This looks like a clear area for future work. It can build on the earlier \gls{ORCID} code.
\Needspace*{18\baselineskip}
    \item Peer-Review Status \& Acceptance Tracking

Because doctoral theses are frequently submitted while papers are in various stages of publication, annotations can dynamically distinguish formal review stages:

\begin{lstlisting}[language={BibTeX}, breaklines=true, basicstyle=\small\ttfamily,  escapechar=|, caption={{Storing Peer-Review Status \& Acceptance Tracking}}, label=lst:biberPeerReviewMetadata]
@article{maguire2026verification,
  author   = {Maguire, Jr., Gerald Q. and Student, Alice},
  entry+an = {peerreview="accepted"; date_accepted="2026-08-15"; preprint="arxiv:2608.01234"},
  title    = {Formal Verification of SmartNIC Microcode},
  journal  = {IEEE Micro},
  year     = {2026},
  placeholder={Placeholder fake reference for use in examples}
}
\end{lstlisting}

\begin{description}
    \item[Use in Document] Bibliography listings can automatically append \textbf{"(Accepted for publication, Aug 2026)"} or \textbf{"(Under peer review)"} without hardcoding temporary text into the \texttt{journal} or \texttt{year} fields.
    
    \item[Automated Guardrails] Build hooks can emit compiler warnings if a student attempts to build a final dissertation (\texttt{\textbackslash finaltrue}) while citations contain \texttt{peerreview="submitted"}.
\end{description}

Before attempting this, more thought is needed --- as this would seem difficult to maintain \textbf{and} many of the manuscripts will not be accepted (except as preprints) before the thesis is finalized and printed. However, one might consider a script-based workflow that checks the relevant \gls{DiVA} entry for updated information.
\end{enumerate}

\subsection{Using annotations for math in titles}
\label{sec:annotationForMathInTitles}

Some bibliographic titles contain \LaTeX{} math expressions. \BibLaTeX{} provides built-in literal fields for sorting and indexing, specifically \texttt{sorttitle} and \texttt{indextitle}, as illustrated in \Cref{lst:biberSorttitle}. These fields decouple sorting and indexing strings from complex typesetting commands.

\begin{listing}[!ht]
\begin{minted}[breaklines, escapeinside=||]{bibtex}
@article{Noz_1975,
  author        = {Noz, Marilyn E. and Maguire, Jr., Gerald Q.},
  title         = {{Calculation of $\bar{E}_{\beta}$, $\Gamma$ and $\Delta_{i}$ for $^{99m}$Tc}},
  sorttitle     = {Calculation of E_beta, Gamma and Delta_i for 99mTc},
  |\dots|
}
\end{minted}
\caption{Using sorttitle}
\label{lst:biberSorttitle}
\end{listing}
\FloatBarrier

While \texttt{sorttitle} normalizes the title into plain ASCII for alphabetical ordering, external ingest pipelines (\eg\ automated bibliography validators, metadata scrapers, or repository systems such as SwePub and \gls{DiVA}) frequently require a clean UTF-8 plain-text representation. We can use \BibLaTeX{} data annotations to store a canonical UTF-8 version of the title, as shown in \Cref{lst:biberMathInTitles}. Note that the leading \enquote{\texttt{=}} (equal sign) instructs \texttt{biber} to treat the annotation value as a literal string rather than a key-value or boolean flag.

\begin{listing}[!ht]
\begin{minted}[breaklines, escapeinside=||]{bibtex}
@article{Noz_1975,
  author        = {Noz, Marilyn E. and Maguire, Jr., Gerald Q.},
  title         = {{Calculation of $\bar{E}_{\beta}$, $\Gamma$ and $\Delta_{i}$ for $^{99m}$Tc}},
  title+an:utf8 = {=Calculation of Ēᵦ, Γ and Δᵢ for ⁹⁹ᵐTc},
  |\dots|
}
\end{minted}
\caption{Using data annotations to represent math in titles}
\label{lst:biberMathInTitles}
\end{listing}
\FloatBarrier

Both approaches can be combined, as shown in \Cref{lst:biberMathInTitlesCombined}. This allows \texttt{sorttitle} to govern internal \texttt{biber} sorting while \texttt{title+an:utf8} supplies external validation pipelines and metadata harvesters with an authoritative UTF-8 string.

\begin{listing}[!ht]
\begin{minted}[breaklines, escapeinside=||]{bibtex}
@article{Noz_1975,
  author        = {Noz, Marilyn E. and Maguire, Jr., Gerald Q.},
  title         = {{Calculation of $\bar{E}_{\beta}$, $\Gamma$ and $\Delta_{i}$ for $^{99m}$Tc}},
  sorttitle     = {Calculation of E_beta, Gamma and Delta_i for 99mTc},
  title+an:utf8 = {=Calculation of Ēᵦ, Γ and Δᵢ for ⁹⁹ᵐTc},
  |\dots|
}
\end{minted}
\caption{Combining data annotations with sorttitle}
\label{lst:biberMathInTitlesCombined}
\end{listing}
\FloatBarrier

\subsubsection{Can annotations be used in conjunction with reference managers}
\label{sec:usingAnnotationsWithReferenceManagers}
Considering the usage \enquote{Thesis Inclusion and Paper Categorization (Compilation Theses)} above, there is a practical problem: 
Currently, no mainstream reference manager supports biber’s data annotations as first-class, structured semantic data.

The landscape divides into tools that parse and preserve the raw fields versus tools that strip them entirely:
\begin{enumerate}
    \item JabRef (Best in Class, but Text-Only)

Because JabRef uses the \fname{.bib} file directly as its internal database (without converting entries into a generic internal schema), it is the most compatible tool:
\begin{description}[labelwidth =\widthof{\textbf{Proprietary Flat Schemas}}, leftmargin = !]
    \item[Preservation] When set to \BibLaTeX{} mode, JabRef will not corrupt or strip `author+an` or `title+an` fields when saving.
    \item[UI Editing] It does not parse the sub-properties (`1:orcid="...";`) into \gls{GUI} checkboxes, creator cards, or linked badges. The entire annotation block appears as an unformatted raw string in an auxiliary field text box.
    \item[Feature Status] Native support for \texttt{biber} data annotations has been discussed in feature issues (\eg Issue \#11505), but it remains an unparsed text field.
\end{description}

    \item Zotero \& Mendeley (Destructive / Incompatible)

Both tools convert bibliographic entries into an internal relational model that assumes a flat \enquote{author} table (\enquote{first\_name}, \enquote{last\_name}, \enquote{role}):
\begin{description}[labelwidth =\widthof{\textbf{Proprietary Flat Schemas}}, leftmargin = !]
    \item[Data Loss on Import] If a \fname{.bib} file decorated with \enquote{author+an = { 1:orcid="..."; }} is imported, the parser discards the annotation entirely because it does not map to any creator field.
    \item[Export Inability] Neither tool can reconstruct \enquote{<field>+an} upon export unless an external layer like Better BibTeX (BBT) parses custom dummy strings stashed in the record's \enquote{Extra} field.
\end{description}

    \item EndNote, Citavi, \& Paperpile
\begin{description}[labelwidth =\widthof{\textbf{Proprietary Flat Schemas}}, leftmargin = !]
    \item[Proprietary Flat Schemas] These tools lack any concept of item/part-level annotation syntax.
    \item[Export Stripping] They adhere to standard legacy BibTeX fields (\texttt{author}, \texttt{title}, \texttt{journal}, \etc). Unknown fields with plus signs (\enquote{+an}) are either sanitized away or converted to malformed entries during export.
\end{description}
\end{enumerate}

\subsubsection{Practical Recommendation}\mbox{}\par
\smallskip

Because reference managers lag behind \texttt{biber}’s advanced data-model features, \textbf{managing annotations outside the \gls{GUI} reference manager is the only stable design pattern}. For a \gls{KTH} doctoral workflow, managing \glspl{ORCID} and \gls{CReDiT} through \fname{publications\_map.json} and external Python/\gls{CI} scripts remains far more robust than relying on a reference manager.

This means:
\begin{enumerate}
    \item \fname{references.bib} as the standard citation repository (which Zotero or JabRef can export safely), and
    \item \fname{publications\_map.json} (managed via the Python/Streamlit wizard and GitHub Actions) as the pipeline that programmatically decorates entries with \enquote{author+an} at build time, ensuring reference managers never accidentally wipe out contributor roles or \glspl{ORCID} during library re-exports.
\end{enumerate}

\subsubsection{A Forward-Looking Architecture: Split-File Overlay}\mbox{}\par
\smallskip

A cleaner architectural path forward is to decouple the reference management lifecycle from semantic metadata authoring by maintaining two distinct bibliographic files sharing identical citation keys:
\begin{enumerate}
\item \fname{references.bib}: The upstream, \gls{GUI}-managed store (exported directly from Zotero, Mendeley, or JabRef) containing canonical bibliographic descriptors (\eg \texttt{author}, \texttt{title}, \texttt{year}, \texttt{doi}).

\item \fname{references-annotations.bib}: A companion overlay containing sparse entries with matching citation keys that declare only the extended data annotations (\eg \texttt{author+an}, \texttt{title+an}, or \texttt{entry+an}).
\end{enumerate}

Existing utilities such as \texttt{bibtool}, \texttt{bibtex-tools}, or \texttt{bibmerge} are designed to detect duplicate citation keys and discard redundant entries. In contrast, an automated Python pre-compilation step (\eg integrated into a GitHub Action or local pre-commit hook) can perform an \emph{additive, field-level unification}. By parsing both files and merging fields that share a matching key into a unified intermediate file (\eg \fname{references-merged.bib}), reference managers can safely overwrite \fname{references.bib} during iterative library exports without clobbering curated author roles or \glspl{ORCID}.

Such an overlay or split-file architecture (\fname{references.bib} alongside \fname{references-annotations.bib}) cleanly separates concerns while providing persistent annotations. As a result, students never have to manually re-annotate files after updating their reference library.

In \texttt{biber} and \BibLaTeX{}, data annotations are completely generic: \emph{any} standard or custom field can take the \enquote{+an} suffix at the item, field, or entry level.

Under the split-file overlay architecture (\fname{references.bib} + \fname{references-annotations.bib}), an additive Python pre-processor can attach granular semantic metadata across several distinct bibliographic dimensions without reference-manager corruption. Examples of these other dimensions include:
\begin{enumerate}
    \item Identifier Verification \& Provenance (\texttt{doi+an}, \texttt{eprint+an}, \texttt{isbn+an})
    \begin{description}[labelwidth =\widthof{\textbf{Application}}, leftmargin = !]
        \item[Context] In automated pipelines, it is crucial to record whether a persistent identifier has been mechanically validated against Crossref, DataCite, MathSciNet, LIBRIS, or another source.
        \Needspace*{6\baselineskip}
        \item[Syntax]
\begin{lstlisting}[language={BibTeX}, breaklines=true, basicstyle=\small\ttfamily,  escapechar=|, caption={{Syntax for Identifier Verification \& Provenance}}, label=lst:biberIdentifierVerificationAndProvenanceMetadata]
doi      = {10.1109/TNET.2025.1234567},
doi+an   = {status="verified"; resolver="crossref"; checked="2026-03-01"},
eprint   = {2501.09876},
eprint+an = {archive="arxiv"; primaryClass="cs.NI"; openaccess="true"}
\end{lstlisting}
        \item[Application] \BibLaTeX{} macros can selectively render a green verification checkmark beside verified \glspl{DOI} or trigger build-time compiler warnings for unverified or dead links.
    \end{description}
    
    \item Open Access \& Licensing Terms (\texttt{url+an}, \texttt{file+an})
        \begin{description}[labelwidth =\widthof{\textbf{Application}}, leftmargin = !]
        \item[Context] Theses submitted to \gls{DiVA} often must comply with funder mandates (\eg Plan S) requiring explicit distinction between closed publisher links, author-accepted manuscripts (AAM), and open-access mirrors.
        \Needspace*{6\baselineskip}
        \item[Syntax]
\begin{lstlisting}[language={BibTeX}, breaklines=true, basicstyle=\small\ttfamily,  escapechar=|, caption={{Syntax for Open Access \& Licensing Terms}}, label=lst:biberOpenAccessAndLicensingTermsMetadata]
url      = {https://doi.org/10.1016/j.comnet.2025.109876},
url+an   = {access="closed"; version="vor"},
url:mirror = {https://urn.kb.se/resolve?urn=urn:nbn:se:kth:diva-345678},
url:mirror+an = {access="gold"; license="cc-by-4.0"; version="aam"}
\end{lstlisting}

        \item[Application] The bibliography driver can conditionally typeset open-access badges, Creative Commons micro-icons, or repository links.
        \end{description}
    
    \item Publication Venue Integrity \& Indexing (\texttt{journal+an}, \texttt{booktitle+an})
        \begin{description}[labelwidth =\widthof{\textbf{Application}}, leftmargin = !]
        \item[Context] Academic reporting increasingly requires tracking journal quartiles (Q1/Q2), Scopus/Web of Science indexing, or conference CORE rankings (A*, A, B).
        \Needspace*{6\baselineskip}
        \item[Syntax]
\begin{lstlisting}[language={BibTeX}, breaklines=true, basicstyle=\small\ttfamily,  escapechar=|, caption={{Syntax for Publication Venue Integrity \& Indexing}}, label=lst:biberPublicationVenueIntegrityAndIndexingMetadata]
journal    = {IEEE Transactions on Networking},
journal+an = {issn="1063-6692"; ranking="core-a"; jcr="q1"; sjr="1.82"},
booktitle  = {Proceedings of ACM SIGCOMM '25},
booktitle+an = {ranking="core-a-star"; acceptance_rate="15.2%"}
\end{lstlisting}

        \item[Application] A student can generate extended bibliographic summaries (or committee review appendices) that automatically display venue selectivity and impact metrics alongside included thesis publications.
        \end{description}
        
    \item Funding Acknowledgment \& Grant Tracking (\texttt{note+an}, \texttt{addendum+an})
        \begin{description}[labelwidth =\widthof{\textbf{Application}}, leftmargin = !]
        \item[Context] Doctoral projects frequently span multiple funding bodies (\eg VR, Vinnova, EU Horizon Europe) where specific publications must trace back to concrete grant IDs.
        \Needspace*{6\baselineskip}
        \item[Syntax]
\begin{lstlisting}[language={BibTeX}, breaklines=true, basicstyle=\small\ttfamily,  escapechar=|, caption={{Syntax for Funding Acknowledgment \& Grant Tracking}}, label=lst:biberFundingAcknowledgmentAndGrantTrackingMetadata]
addendum   = {Supported by the Swedish Research Council},
addendum+an = {funder="VR"; grant="2023-04567"; project="6G-SEC"}
\end{lstlisting}

        \item[Application] Citation drivers can automatically inject formal grant disclosures into the printed thesis bibliography or export structured funding metadata into \gls{DiVA}/\gls{MODS} exports.
        \end{description}
    \item Editorial Curation \& Retraction Watches (\texttt{entry+an})
        \begin{description}[labelwidth =\widthof{\textbf{Application}}, leftmargin = !]
        \item[Context] Long-lived literature repositories occasionally contain retracted papers, expressions of concern, or \foreignlanguage{latin}{errata}.
        \Needspace*{6\baselineskip}
        \item[Syntax]
\begin{lstlisting}[language={BibTeX}, breaklines=true, basicstyle=\small\ttfamily,  escapechar=|, caption={{Syntax for Editorial Curation \& Retraction Watches}}, label=lst:biberEditorialCurationAndRetractionWatchesMetadata]
entry+an = {
  retraction_status="clean";
  verified_retractionwatch="2026-08-10";
  curation_notes="Benchmark baseline verified against source code"
}
\end{lstlisting}

        \item[Application] Pre-flight \gls{CI} actions can check Crossref's \enquote{is-retracted} \gls{API} (or other source of retractions) and inject warnings into the \gls{PDF} during compilation if an author inadvertently cites retracted work.
        \end{description}
        
    \item Typographic \& Pagination Overrides (\texttt{pages+an}, \texttt{edition+an})
        \begin{description}[labelwidth =\widthof{\textbf{Application}}, leftmargin = !]
        \item[Context] Older conference proceedings, electronic article numbering (e-locators), or bilingual editions often require localized presentation strings without modifying the raw page numbering data.
        \Needspace*{6\baselineskip}
        \item[Syntax]
\begin{lstlisting}[language={BibTeX}, breaklines=true, basicstyle=\small\ttfamily,  escapechar=|, caption={{Syntax for Typographic \& Pagination Overrides}}, label=lst:biberTypographicAndPaginationOverridesMetadata]
pages    = {45--60},
pages+an = {type="article-number"; prefix="art: "},
edition  = {2},
edition+an = {swedish="2:a uppl."; english="2nd ed."}
\end{lstlisting}

        \item[Application] Custom formatting drivers can switch between language-sensitive labels without maintaining multiple \fname{.bib} variants.
        \end{description}
\end{enumerate}

Because reference managers routinely discard or corrupt non-standard fields, isolating these annotations inside \fname{references-annotations.bib} and merging them via Python guarantees that extended provenance, venue, and licensing data remain permanently intact.

\subsection{Contrasting annotations with current workflows}
\label{sec:contrastingAnnotationsWithCurrentWorkflows}
The current design approaches publication management and thesis inclusion from an architectural standpoint that is fundamentally different from a purely in-\TeX{} data-annotation model.

Rather than overloading \fname{.bib} entries with document-assembly directives, the current workflows decouple bibliographic citation data from administrative pipeline state:
\begin{description}[style=nextline, labelwidth = {1cm}, leftmargin = !]
    \item[External Source of Truth (\fname{publications\_map.json})]
In the current workflow, inclusion (\texttt{status: "included"}), label assignments (\texttt{paper:A}, \texttt{techreport:A}), tab order (\texttt{tab\_index}), physical reprint paths (\texttt{file\_path}), page slices (\texttt{pdf\_pages}), scaling factors, and reproduction copyright notices (\texttt{permission\_text}) live entirely inside \fname{publications\_map.json}. This separates immutable bibliographic facts (which belong in \fname{references.bib}) from the mutable, student-curated decisions about how physical \gls{PDF} reprints are laid out in a compilation thesis.


    \item[Two-Way Upstream Sync (\gls{DiVA} Discovery)]
The automated workflow pulls publication metadata directly from \gls{DiVA} via the student's KTHID through GitHub Actions. In a pure in-\TeX{} annotation workflow, the candidate must hand-annotate raw \fname{.bib} files. In the current pipeline, the GitHub Action generates \fname{lib/publications\_generated.tex} and \fname{lib/publications\_dividers\_generated.tex} dynamically from \gls{DiVA} and \fname{publications\_map.json}, removing the manual editing step entirely.

\smallskip\noindent
\item[CI/CD Build-Time Validation]
Because paper inclusion is tracked in \gls{JSON}, the GitHub Actions can validate physical \gls{PDF} assets \textbf{before} typesetting:

\smallskip
\noindent\makebox[\linewidth][l]{%
  \hspace*{-\leftmargin}%
  \footnotesize
  \begin{tabular}{@{}ll@{}}
    \texttt{"status": "included", "pdf\_downloaded": false} & $\longrightarrow$ Emits visual missing-file warning page\\
    \texttt{"status": "included", "pdf\_downloaded": true}  & $\longrightarrow$ Injects divider tab and binds reprint PDF
  \end{tabular}%
}
\smallskip

Testing file existence, checking scaling, and enforcing page intervals are trivial in Python, whereas executing robust filesystem checks and multi-page \gls{PDF} inclusions dynamically inside macro-level \texttt{biblatex} hooks is notoriously fragile.


    \item[Granular Multi-Author \& Collaboration CReDiT]
Handling of massive collaborations (such as, the BNL Multi-Particle Spectrometer example) via \texttt{specific\_contributors} demonstrates why the wizard approach (\texttt{CReDiT\_Matrix\_Wizard.py}) is cleaner than annotating \fname{.bib} files. In large multi-author experiments, attempting to attach 14 \gls{CReDiT} part annotations to individual authors inside a standard BibTeX string quickly becomes unwieldy. Generating a standalone \fname{lib/thesis\_contributions\_generated.tex} file keeps the bibliography clean while producing an explicit, readable contribution matrix for the defense committee.
\end{description}

The existing pipeline positions the \fname{.bib} file as a traditional citation store, delegating document assembly, inclusion filtering, DiVA harvesting, and \gls{PDF} insertion to external Python and \gls{CI} automation where it is most maintainable.

\subsection{Support for Mega-Collaborations}
\label{sec:biberMegaCollaborations}
To prevent \texttt{biber} and \LaTeX{} from crashing on mega-collaborations while still providing granular \gls{CReDiT} attribution for a student, you can use \textbf{sparse target indexing}, (\ie you only need to add the relevant additional author details):

\Needspace*{24\baselineskip}
\begin{lstlisting}[language={BibTeX}, breaklines=true, basicstyle=\small\ttfamily,  escapechar=|, caption={{Example of a specific contributor to a large collaboration}}, label=lst:biberMegaCollaborationMetadata]
@techreport{bnlMPS1975,
  author     = {Lindenbaum, S. J. and {MPS Collaboration}},
  author+an  = {
    3:specific-contributor="Maguire Jr., Gerald Q.";
    3:software="lead";
    3:visualization="lead";
  },
  title      = {The {Brookhaven} {Multi-Particle} {Spectrometer} ({MPS}) Facility},
  year       = {1975},
  number     = {BNL-20202},
  url        = {https://www.osti.gov/servlets/purl/4162299},
  note       = {An invited paper presented at the IV International Symposium on NN Interaction, Syracuse University, May 2-4, 1975.},
  entry+an   = {
    diva_id="bnlMPS1975";
    contribution_note="Summer student working on PDP-10 software and graphical display of data.";
    has_specific_contributors="true";
  },
  institution = {Brookhaven National Laboratory},
  type       = {Report},
}
\end{lstlisting}

\subsection{Integrating annotations with current workflows}
\label{sec:integratingAnnotationsWithCurrentWorkflows}

The challenge with handling entities like the \textbf{MPS Collaboration} in standard \fname{.bib} files is that the individual student or researcher (\eg Gerald Q. Maguire Jr.) is not listed as an indexed name item in the \texttt{author} field (which contains only \enquote{Lindenbaum, S. J.} and \enquote{\{MPS Collaboration\}}). Consequently, biber’s item-indexed annotations (\texttt{author+an = \{ 3:software="supporting"\}}) cannot bind to an author index (in this example \enquote{3}) that does not exist in the primary name list (or might only appear far down the list, \ie beyond the range of the authors processed).

Under biber, this can be solved cleanly using one of two annotation architectures:

\begin{enumerate}
    \item Standalone Contributor Data Annotation (\texttt{contributors+an})

Rather than attempting to attach metadata to the primary \texttt{author} field, declare an annotated dummy or user field (such as \texttt{contributors+an}). Biber allows data annotations on arbitrary fields, even when the underlying base field is sparse or virtual.

In your \fname{references-annotations.bib} overlay:

\begin{lstlisting}[language={BibTeX}, breaklines=true, basicstyle=\small\ttfamily,  escapechar=|, caption={{Example of non-named author collaborator inclusion}}, label=lst:biberPatternOneCollaborMetadata]
@techreport{bnlMPS1975,
  % Entry annotation to flag a collaborator inclusion
  entry+an = {
    has_specific_contributors = "true";
    collaboration_name        = "MPS Collaboration";
    contribution_note         = "Summer student working on PDP-10 software and graphical display of data."
  },
  % Granular CReDiT annotations for contributors not named in the author field:
  contributors+an = {
    1:name="Maguire Jr., Gerald Q.";
    1:role="software";
    1:role="visualization";
    1:affiliation="Brookhaven National Laboratory"
  }
}
\end{lstlisting}

\Needspace*{13\baselineskip}
\textbf{Inside LaTeX / BibLaTeX:}

You can query this with:
\begin{lstlisting}[language={[LaTeX]TeX}, breaklines=true, basicstyle=\small\ttfamily,  escapechar=|, caption={{Standalone Contributor Data Annotation code}}, label=lst:biberPatternOneCode]
\hasentryannotation{has_specific_contributors}{%
  \par\smallskip\noindent\textbf{Collaboration Contributor Disclosure:}%
  \par\hspace*{1.5em}\getentryannotation{contribution_note}%
  \par\hspace*{1.5em}\getpartannotation[contributors][default][1]{name}: %
    \textit{\getpartannotation[contributors][default][1]{role}}%
}{}
\end{lstlisting}

    \item Part Annotation on Collaboration Author Index (Sub-part Mapping)

If the collaboration is listed in the \texttt{author} field as a collective corporate entity (for instance, index 2 in \texttt{Lindenbaum, S. J.} and \texttt{\{MPS Collaboration\}}), you can attach the individual contributor directly to that author's index using structured part keys:
\Needspace*{15\baselineskip}
\begin{lstlisting}[language={BibTeX}, breaklines=true, basicstyle=\small\ttfamily,  escapechar=|, caption={{Syntax for Part Annotation on Collaboration Author Index (Sub-part Mapping)}}, label=lst:biberPatternTwoCollaborPartMetadata]
@techreport{bnlMPS1975,
  author    = {Lindenbaum, S. J. and {MPS Collaboration}},
  author+an = {
    % Roles for Professor Lindenbaum
    1:conceptualization="lead";
    1:supervision="lead";
    % Specific contributor anchored to Author 2 (the Collaboration)
    2:specific-contributor="Maguire Jr., Gerald Q.";
    2:software="supporting";
    2:visualization="supporting";
    2:note="Summer student: PDP-10 assembly-to-FORTRAN conversion & track display"
  }
}
\end{lstlisting}

\Needspace*{14\baselineskip}
\textbf{How the In-TeX Driver Handles It:}

In your existing \texttt{creditnames} name formatting driver, when the driver processes Author 2 (\texttt{MPS Collaboration}), you can add an interceptor:
\begin{lstlisting}[language={[LaTeX]TeX}, breaklines=true, basicstyle=\small\ttfamily,  escapechar=|, caption={{Code for handling specific-contributor metadata}}, label=lst:biberPatternTwoCollaborPartMetadataCode]
\newcommand*{\checkspecificcontributors}[1]{%
  % #1 = current author index
  \haspartannotation[author][default][#1]{specific-contributor}{%
    \par\hspace*{1.5em}\textbf{Specific Contributor:} %
      \getpartannotation[author][default][#1]{specific-contributor}%
    \haspartannotation[author][default][#1]{note}{%
      \par\hspace*{3em}\textit{\getpartannotation[author][default][#1]{note}}%
    }{}%
  }{}%
}
\end{lstlisting}
\end{enumerate} 




\Needspace*{18\baselineskip}
\textbf{Python Pipeline Synthesis (Connecting \texttt{CReDiT\_Matrix\_Wizard.py})}

In the Streamlit wizard and GitHub Actions workflow:
\begin{enumerate}
    \item \fname{publications\_map.json} remains the interactive store holding\linebreak[4] \texttt{specific\_contributors} and \texttt{contribution\_note}.
    
    \item When the wizard generates or updates \fname{references-annotations.bib}, the Python script inspects \texttt{entry.get("specific\_contributors")}.
    
    \item If \texttt{specific\_contributors} contains entries, Python matches the named collaboration in the author list and writes the corresponding \texttt{2:specific-contributor="\ldots"} and role slugs directly into the generated overlay entry.
\end{enumerate}

This preserves standard bibliography listings for external readers while enabling the defense committee to inspect candidate contributions to large institutional consortia.

It should be possible to mine the \fname{publications\_map.json} file generated by \fname{CReDiT\_Matrix\_Wizard.py} to produce an initial \fname{references-annotations.bib} file with the \gls{CReDiT} annotations for authors.

\subsection{Scaling with the number of publications a student has in DiVA}
The \fname{publications\_map.json} file has the same number of entries as the number of publications (other than the current thesis) that an author already has in \gls{DiVA}. How large is this likely to be? \Cref{tab:histogramPriorPublications} shows how many students (from those who defended in 2025) have a given number of other publications in \gls{DiVA} (\ie not counting any third-cycle theses).

From this data, we see that the distribution is heavily right-skewed by outliers with 46 and 405 publications in \gls{DiVA}. Looking at these, we see that the student with 405 publications in \gls{DiVA} worked in the ATLAS consortium (thesis title: \textit{Divide and Conquer: A Fine-Grained Look at the Higgs Boson Cross Section in WW Decays Using the ATLAS Experiment}). This means that if they were to identify their contributions to specific papers, the student would have to use the method described in \Cref{sec:biberMegaCollaborations}. However, in this case, the student wrote a monograph and not a compilation-style thesis.
In the case of the student with 46 publications in Electromagnetic Engineering (thesis title: \textit{Modeling and design of lenses combined with array antennas in the \mbox{millimetre-wave} regime}), the student chose to prune, and they included seven publications in their compilation thesis and cited five others that were related.

Of the $N = 310$ theses (comprising both doctoral and licentiate degrees), the median is 7 prior publications, with an \gls{IQR} spanning 4 to 10. Less than \qty{2}{\percent} of students have 36 or more publications in \gls{DiVA}, and only \qty{4.2}{\percent} have 25 or more. Therefore, the manual effort required to designate which publications are to be included in a compilation thesis\footnote{Unfortunately, \gls{DiVA} does not provide an automated mechanism to query constituent paper linkages for third-cycle theses in aggregate. Consequently, determining the exact number of publications formally incorporated into each compilation thesis requires inspecting each individual record or \gls{PDF}. \Cref{sec:includedPublications2025} examines manually collected data for the year 2025.} and specify their tab order and logical labels remains feasible. Furthermore, because authors include only a curated subset of their total output, subsequent role curation in \texttt{CReDiT\_Matrix\_Wizard.py} remains a modest-sized task.

\begin{table}[!htp]
    \caption{Histogram of the number of authors with a given number of publications in \gls{DiVA} for third-cycle students who presented in 2025.}
    \label{tab:histogramPriorPublications}
    \centering
    \begin{tabular}{S[table-format=4.0] S[table-format=4.0]}
    \toprule
      {\textbf{Number of other publications in \gls{DiVA}}} & {\textbf{Number of authors}}\\
    \midrule
1&	6\\
2&	15\\
3&	24\\
4&	31\\
5&	28\\
6&	22\\
7&	31\\
8&	27\\
9&	14\\
10&	17\\
11&	15\\
12&	14\\
13&	12\\
14&	11\\
15&	9\\
16&	5\\
17&	5\\
18&	2\\
19&	3\\
21&	1\\
22&	1\\
24&	4\\
25&	5\\
26&	1\\
27&	1\\
28&	1\\
34&	1\\
36&	2\\
46&	1\\
405&	1\\
\bottomrule
\end{tabular}
\end{table}
\FloatBarrier

\subsubsection{Summary of Included Publications in Compilation theses (2025 Cohort)}
\label{sec:includedPublications2025}
Out of the 310 total third-cycle theses presented in 2025, exactly 298 (\qty{96.1}{\percent}) were compilation-style theses with constituent publications, leaving 12 monographs. The number of included publications in these compilation-style theses from 2025 was manually collected from \gls{DiVA} ($N = 298$, comprising 271 doctoral and 27 licentiate theses). \Cref{tab:histogramIncludedPublications2025} gives the distribution of the number of included publications, while the summary below outlines key statistical measures:

\begin{description}[style=nextline, labelwidth=\widthof{\textbf{Combined Cohort ($N = 298$)}}, leftmargin=!]
    \item[Combined Cohort ($N = 298$)]
        Mean: 4.89; Standard Deviation ($s$, sample/pop): 1.66 / 1.65; Median: 5.0; Mode: 4 (110 theses); Range: 2 to 12; $Q_1 = 4.0$, $Q_3 = 6.0$, $\text{IQR} = 2.0$.
    \item[Doctoral Theses ($N = 271$)]
        Mean: 5.13; Standard Deviation ($s$, sample/pop): 1.53 / 1.53; Median: 5.0; Mode: 4 (109 theses); Range: 2 to 12; $Q_1 = 4.0$, $Q_3 = 6.0$, $\text{IQR} = 2.0$.
    \item[Licentiate Theses ($N = 27$)]
        Mean: 2.48; Standard Deviation ($s$, sample/pop): 0.58 / 0.57; Median: 2.0; Mode: 2 (15 theses); Range: 2 to 4; $Q_1 = 2.0$, $Q_3 = 3.0$, $\text{IQR} = 1.0$.
\end{description}

\begin{table}[!htp]
    \caption{Distribution of included publications in third-cycle compilation theses presented in 2025.}
    \label{tab:histogramIncludedPublications2025}
    \centering
    \begin{tabular}{S[table-format=2.0] S[table-format=3.0] S[table-format=2.2] S[table-format=2.0] S[table-format=2.2] S[table-format=3.0] S[table-format=2.2]}
    \toprule
      & \multicolumn{2}{c}{\textbf{Doctoral ($N=271$)}} & \multicolumn{2}{c}{\textbf{Licentiate ($N=27$)}} & \multicolumn{2}{c}{\textbf{Combined ($N=298$)}} \\
      \cmidrule(lr){2-3} \cmidrule(lr){4-5} \cmidrule(lr){6-7}
      {\textbf{Publications}} & {\textbf{Freq.}} & {\textbf{\%}} & {\textbf{Freq.}} & {\textbf{\%}} & {\textbf{Freq.}} & {\textbf{\%}} \\
    \midrule
      2  &   4 &  1.48 & 15   & 55.56 &  19 &  6.38 \\
      3  &   6 &  2.21 & 11   & 40.74 &  17 &  5.70 \\
      4  & 109 & 40.22 &  1   &  3.70 & 110 & 36.91 \\
      5  &  66 & 24.35 & {--} & {--}  &  66 & 22.15 \\
      6  &  43 & 15.87 & {--} & {--}  &  43 & 14.43 \\
      7  &  24 &  8.86 & {--} & {--}  &  24 &  8.05 \\
      8  &   9 &  3.32 & {--} & {--}  &   9 &  3.02 \\
      9  &   5 &  1.85 & {--} & {--}  &   5 &  1.68 \\
     10  &   2 &  0.74 & {--} & {--}  &   2 &  0.67 \\
     11  &   2 &  0.74 & {--} & {--}  &   2 &  0.67 \\
     12  &   1 &  0.37 & {--} & {--}  &   1 &  0.34 \\
    \bottomrule
    \end{tabular}
\end{table}
\FloatBarrier



\subsubsection{Key Analytical Takeaways With Respect to Included Publications}
\begin{description}[style=nextline, labelwidth=\widthof{\textbf{CReDiT Tooling Workload}}, leftmargin=!]
\item[Compact Scaling Envelope] Over \qty{80.4}{\percent} of doctoral compilation theses include between 4 and 6 papers, with the modal value being exactly 4.

\item[Licentiate Norm] Over \qty{96.3}{\percent} of licentiates incorporate either 2 or 3 papers.

\item[CReDiT Tooling Workload] Even for the single extreme doctoral outlier with 12 included works, populating \fname{publications\_map.json} and managing contributor attributions in\linebreak[4] \texttt{CReDiT\_Matrix\_Wizard.py} involves an input matrix of modest dimensions. For most doctoral candidates ($Q_1 = 4$, $Q_3 = 6$, $\mathrm{IQR} = 2$), the manual CReDiT assignment task requires processing only 4 to 6 rows.
\end{description}


\section{Recommendations for generating fake \BibLaTeX{} entries}
\label{sec:generatingFakeBibLaTeXEntries}

\gls{ORCID} iDs are a designated subset of \glspl{ISNI}\footnote{See \url{https://isni.org/} and \gls{ISO} 27729~\cite{iso27729:2024} for more information.}, specifically within the range \texttt{0000-0001-5000-0007} to \texttt{0000-0003-5000-0001}. Values entirely outside this assignment block --- such as those starting with all zeroes (\texttt{0000-0000-XXXX-XXXX}) --- do not represent valid \gls{ORCID} registry allocations. Additionally, \gls{ORCID} iDs use a MOD 11-2 checksum as the final character; thus, any 16-digit string with a failed checksum is invalid and will be rejected by any schema or \gls{API} validator. Because valid \glspl{ISNI} exist outside the designated \gls{ORCID} block, arbitrary non-ORCID \glspl{ISNI} should still be avoided to prevent collisions with real registered entities in the global \gls{ISNI} database; test entries should strictly employ mathematically invalid checksums or zero-filled blocks (\eg \texttt{0000-0000-0000-000X}).

The International DOI Foundation (IDF) has reserved a designated \gls{DOI} prefix specifically for testing, documentation, and development, and a number of other \gls{DOI} prefixes are used/reserved for similar purposes. These are:
\begin{itemize}
    \item \texttt{10.1000} (The Demonstration \& Documentation Prefix):
The \textit{DOI Handbook} explicitly assigns the prefix \texttt{10.1000} for examples, demonstration systems, and documentation (similar to example domain names like \enquote{\texttt{example.com}}). Identifiers using this prefix --- such as \texttt{10.1000/182} --- are reserved for instructional and testing purposes and will not resolve to real production literature.

    \item DataCite Test Environment (\texttt{10.5072}):
DataCite reserves prefix \texttt{10.5072} strictly for sandbox environments and test accounts. \glspl{DOI} registered under \texttt{10.5072} are test records used during API integration and are never routed to the production registry.

    \item Crossref Test Prefix (\texttt{10.5555}):
Crossref provides test accounts with prefixes such as \texttt{10.5555} (and similar designated sandbox prefixes) for verifying \gls{XML} deposits and validation pipelines without creating permanent scholarly records.

    \item Reserved Registrant Code \texttt{10.0000}:
The \texttt{10.0000} space is unassigned to any \gls{DOI} Registration Agency (RA) and is invalid for production publication metadata.
\end{itemize}

For mock datasets, documentation examples, or pipeline test suites, \texttt{10.1000} is the standard, RFC/handbook-compliant choice. Use this prefix for examples.

\subsection{Authors, Titles, and other Metadata Fields}

When constructing synthetic entries for documentation, testing, or pipeline verification, all accompanying fields must be constructed so that they cannot be conflated with genuine scholarly assets:

\begin{description}[style=nextline, labelwidth=\widthof{\textbf{Domains \& URIs}}, leftmargin=!]
    \item[Explicit Metadata Markers (\texttt{placeholder})]
    When real names, institutional affiliations, or historical works are used in illustrative examples, the entry should explicitly declare:
\begin{lstlisting}[language={BibTeX}, basicstyle=\small\ttfamily]
placeholder = {Placeholder fake reference for use in examples},
\end{lstlisting}
    This provides a machine-readable flag that downstream \gls{CI} verification scripts (\eg \texttt{bib\_cleanup.py}) and pre-flight linters can inspect to ensure synthetic records are not inadvertently retained in final thesis submissions or uploaded to production repositories like \gls{DiVA}.

    \item[Authors \& Consortia]
    Avoid realistic common names that could collide with actual researchers in databases such as Scopus, Web of Science, or DiVA. Prefer conventional metasyntactic names (\eg \enquote{Alice Example} and \enquote{Bob Sample}) or clearly designated institutional placeholders (\eg \enquote{Researcher, Mock A.} or \texttt{\{\{The Example Collaboration\}\}}).

    \item[Titles \& Publication Venues]
    Titles should explicitly include the terms \enquote{Example}, \enquote{Mock}, or \enquote{Demonstration} (\eg \textit{A Synthetic Benchmark for Example Systems}). Journals and book titles should use explicitly fictional titles (\eg \textit{Journal of Illustrative Studies}, \textit{Proceedings of the Demonstration Conference}) to avoid collisions with predatory or newly minted real journals.

    \item[ISBNs \& ISSNs]
    Follow \gls{ISO} 2108 recommendations for example \glspl{ISBN}~\cite{iso2108:2017}. Use the standard demonstration ISBN-13 (\texttt{978-3-16-148410-0}) or an unassigned/failed-checksum sequence (\eg \texttt{978-0-00-000000-0}). For serial publications, the dummy \gls{ISSN} \texttt{0000-0000} is unallocated by the \gls{ISSN} International Centre.

    \item[Domains \& URIs]
    Under \gls{RFC} 2606~\cite{rfc2606} and its update \gls{RFC} 6761~\cite{rfc6761}, only the reserved second-level domains (\texttt{example.com}, \texttt{example.org}, \texttt{example.edu}) or the reserved top-level domains (\texttt{.test}, \texttt{.example}, \texttt{.invalid}) should be used for \glspl{URL} and repository mirrors. For national repository handles, use zero-padded identifiers (\eg \texttt{urn:nbn:se:kth:diva-000000}).
\end{description}

\section{Overleaf Automation}
\label{sec:verleafAutomation}

\Cref{{sec:AutomatingPlaceholdCheck}} examines how a script could be used to check for placeholders. \Cref{sec:furtherAutomation} generalizes this approach to copy files from the compilation container to the Overleaf project, so that artifacts can become persistent and potentially be stored in a GitHub repository.

\subsection{Automating the Placeholder check}
\label{sec:AutomatingPlaceholdCheck}

Overleaf's internal build pipeline is fundamentally stateless with respect to the git tree, as the compiler container writes auxiliary files, \enquote{.log}, and \enquote{.pdf} to an ephemeral build cache outside the git repository directory. The web \gls{UI} exposes these via internal project endpoints that are not exposed in standard git syncs.

Breaking the \gls{SaaS} walled gardens via the UNIX philosophy: \textbf{If a browser button triggers an action, there is an underlying HTTP endpoint or a DOM event that can be automated.}


There are two primary approaches to automating this retrieval without manually clicking through the \gls{GUI}:
\begin{enumerate}[leftmargin=*, label=\textbf{Method \arabic*}, ref=Method \arabic*]
\item Use a script to compile the file, then fetch the output.log file and filter it

If we assume that you are logged in and accessing your Overleaf project using Chrome as your browser. You should set the environment variable \texttt{OVERLEAF\_PROJECT\_ID} to be the identifier of your Overleaf project, \ie the string that follows \enquote{project/}. \Cref{lst:accessBuildLogDirectly} shows a Python program that accesses Overleaf (via its \gls{CLSI}), authenticating with your standard session cookie (\texttt{overleaf\_session2}).  \Cref{lst:bashForRemainingPlaceholders} shows the output of the program when accessing this \LaTeX{} file. The hardest part of developing this program was programmatically discovering all MongoDB 24-hex ObjectIds. The program can take one of three arguments (see \texttt{resolve\_file\_target}): \first the 24-hexadecimal ID, \Second the exact path name, or\third part of a path name. If the program finds unresolved placeholders, it returns \num{2} so it can be used in a \gls{CI} process.

\begin{lstlisting}[language={Python}, caption={Use Python to access the build log},label=lst:accessBuildLogDirectly]
#!/usr/bin/env python3
"""
fetch_overleaf_check_for_unresolved_placeholders4.py

Automates Overleaf log retrieval and audits for unresolved BibLaTeX placeholder
references. Programmatically discovers all MongoDB 24-hex ObjectIds directly
from Overleaf's /updates and /entities endpoints without requiring manual
JSON mapping files or browser developer window scripts.
"""

import os
import re
import sys
import uuid
import requests

try:
    import browser_cookie3
except ImportError:
    browser_cookie3 = None

TARGET_BANNER = "UNRESOLVED PLACEHOLDER REFERENCES DETECTED IN THIS BUILD:"
PLACEHOLDER_MARKER = "Unresolved placeholder reference"

FallBack_Project_ID="000000000000000000000000" # replace with your project ID
FallBack_CLSIserver="clsi-cache-zone-b-prod-3" # replace with your CLSI server


def get_authenticated_session(project_id: str) -> tuple[requests.Session, str]:
    session = requests.Session()
    session.headers.update({
        "User-Agent": (
            "Mozilla/5.0 (X11; Linux x86_64) AppleWebKit/537.36 "
            "(KHTML, like Gecko) Chrome/154.0.0.0 Safari/537.36"
        ),
        "Referer": f"https://www.overleaf.com/project/{project_id}",
        "Origin": "https://www.overleaf.com",
        "Accept": "application/json, text/plain, */*",
    })

    cookie_val = os.environ.get("OVERLEAF_SESSION_COOKIE")

    if not cookie_val and browser_cookie3:
        try:
            cj = browser_cookie3.chrome(domain_name=".overleaf.com")
            session.cookies.update(cj)
            for c in cj:
                if c.name == "overleaf_session2":
                    cookie_val = c.value
                    break
        except Exception as e:
            sys.stderr.write(f"Note: browser_cookie3: {e}\n")

    if cookie_val:
        if "overleaf_session2=" in cookie_val:
            cookie_val = cookie_val.split("overleaf_session2=")[1].split(";")[0].strip()
        session.cookies.set("overleaf_session2", cookie_val, domain=".overleaf.com", path="/")

    res = session.get(f"https://www.overleaf.com/project/{project_id}")
    if res.status_code == 403 or "login" in res.url:
        sys.stderr.write("Authentication failed. Ensure you are logged into Overleaf in Chrome.\n")
        sys.exit(1)

    csrf_token = None
    for pat in [
        r'window\.csrfToken\s*=\s*["\']([^"\']+)["\']',
        r'<meta\s+name=["\']ol-csrfToken["\']\s+content=["\']([^"\']+)["\']',
        r'"_csrf"\s*:\s*["\']([^"\']+)["\']',
    ]:
        m = re.search(pat, res.text)
        if m:
            csrf_token = m.group(1)
            break

    if not csrf_token:
        sys.stderr.write("Failed to retrieve CSRF token from project page.\n")
        sys.exit(1)

    return session, csrf_token


def fetch_project_doc_mapping(session: requests.Session, project_id: str, csrf_token: str) -> dict[str, str]:
    """
    Queries Overleaf's /updates endpoint to dynamically extract the mapping between
    file relative paths and their 24-hex MongoDB ObjectIds.
    """
    mapping = {}
    headers = {
        "X-Csrf-Token": csrf_token,
        "X-Requested-With": "XMLHttpRequest",
        "Accept": "application/json",
    }

    try:
        r = session.get(f"https://www.overleaf.com/project/{project_id}/updates", headers=headers, timeout=15)
        if r.status_code == 200:
            data = r.json()
            updates = data.get("updates", [])
            for u in updates:
                # Overleaf update records store pathnames and document/blob metadata
                pathnames = u.get("pathnames", [])
                
                # Check for explicit docs/files lists inside updates
                docs = u.get("docs", []) or u.get("meta", {}).get("docs", [])
                for d in docs:
                    d_id = d.get("id") or d.get("_id")
                    d_path = d.get("pathname") or d.get("path")
                    if d_id and d_path:
                        clean_p = d_path.lstrip("/")
                        mapping[clean_p] = d_id
                        mapping[os.path.basename(clean_p)] = d_id

                # Inspect patches/ops where doc IDs are paired with paths
                patches = u.get("patches", [])
                for p in patches:
                    d_id = p.get("doc") or p.get("doc_id")
                    d_path = p.get("pathname") or p.get("path")
                    if d_id and d_path:
                        clean_p = d_path.lstrip("/")
                        mapping[clean_p] = d_id
                        mapping[os.path.basename(clean_p)] = d_id

                # If updates associate a single doc_id with pathnames
                u_doc_id = u.get("doc") or u.get("doc_id") or u.get("id")
                if u_doc_id and re.fullmatch(r"[a-f0-9]{24}", str(u_doc_id)):
                    for path in pathnames:
                        clean_p = path.lstrip("/")
                        mapping[clean_p] = u_doc_id
                        mapping[os.path.basename(clean_p)] = u_doc_id

            # Fallback regex scan over the payload for any "pathname":"...", "id":"..." structures
            if not mapping:
                for match in re.finditer(
                    r'(?:pathname|path)["\']?\s*:\s*["\']([^"\']+)["\'].*?(?:id|_id|doc)["\']?\s*:\s*["\']([a-f0-9]{24})["\']',
                    r.text
                ):
                    p, did = match.group(1).lstrip("/"), match.group(2)
                    mapping[p] = did
                    mapping[os.path.basename(p)] = did
    except Exception as e:
        sys.stderr.write(f"Warning: Failed to fetch updates mapping: {e}\n")

    return mapping


def fetch_project_paths(session: requests.Session, project_id: str, csrf_token: str) -> list[str]:
    """Retrieves canonical paths from the /entities endpoint."""
    headers = {
        "X-Csrf-Token": csrf_token,
        "X-Requested-With": "XMLHttpRequest",
        "Accept": "application/json",
    }
    try:
        r = session.get(
            f"https://www.overleaf.com/project/{project_id}/entities?projectId={project_id}",
            headers=headers,
            timeout=10
        )
        if r.status_code == 200:
            data = r.json()
            entities = data.get("entities", [])
            return [e.get("path", "").lstrip("/") for e in entities if "path" in e]
    except Exception:
        pass
    return []


def resolve_file_target(
    arg: str | None,
    mapping: dict[str, str],
    known_paths: list[str],
) -> tuple[str | None, str | None]:
    """Resolves an input argument to (root_doc_id, root_resource_path)."""
    if not arg:
        return None, None

    clean = arg.strip().lstrip("./")

    # Invert mapping to support doc_id -> canonical path
    id_to_path = {}
    for p, did in mapping.items():
        if "/" in p:
            id_to_path[did] = p
    for p, did in mapping.items():
        id_to_path.setdefault(did, p)

    # 1. Direct 24-character hexadecimal MongoDB ObjectId
    if re.fullmatch(r"[a-f0-9]{24}", clean):
        doc_id = clean
        path = id_to_path.get(doc_id)
        if not path:
            path = next((p for p in known_paths if os.path.basename(p) == clean), None)
        return doc_id, path

    # 2. Exact match in mapping
    base = os.path.basename(clean)
    doc_id = mapping.get(clean) or mapping.get(base)

    # 3. Substring / fuzzy match across mapped paths and known entities
    if not doc_id:
        candidates = {k: v for k, v in mapping.items() if base.lower() in k.lower()}
        unique_ids = set(candidates.values())
        if len(unique_ids) == 1:
            doc_id = next(iter(unique_ids))
            clean = max([k for k in candidates if candidates[k] == doc_id], key=len)
        elif len(unique_ids) > 1:
            sys.stderr.write(f"Ambiguous match for '{arg}':\n")
            for k in sorted(candidates.keys()):
                sys.stderr.write(f"  * {k}\n")
            sys.exit(1)

    # 4. Resolve relative canonical directory path
    full_path = id_to_path.get(doc_id)
    if not full_path:
        for p in known_paths:
            if os.path.basename(p) == base or p == clean:
                full_path = p
                break

    if not full_path:
        full_path = clean if "/" in clean else base

    return doc_id, full_path


def fetch_latest_log(
    session: requests.Session,
    project_id: str,
    csrf_token: str,
    root_doc_id: str | None = None,
    root_resource_path: str | None = None,
) -> str:
    compile_url = f"https://www.overleaf.com/project/{project_id}/compile?enable_pdf_caching=true"
    headers = {
        "Content-Type": "application/json",
        "X-Csrf-Token": csrf_token,
    }

    payload = {
        "draft": False,
        "png2pdf": True,
        "incrementalCompilesEnabled": True,
        "stopOnFirstError": False,
        "editorId": str(uuid.uuid4()),
    }

    # Pass the resolved MongoDB ObjectId and canonical resource path
    if root_doc_id:
        payload["rootDoc_id"] = root_doc_id
    if root_resource_path:
        payload["rootResourcePath"] = root_resource_path

    # Only request a silent check when building the default root document
    if not root_doc_id and not root_resource_path:
        payload["check"] = "silent"

    compile_res = session.post(compile_url, headers=headers, json=payload)
    if compile_res.status_code != 200:
        sys.stderr.write(f"POST compile failed: HTTP {compile_res.status_code}\n")
        sys.stderr.write(compile_res.text[:300] + "\n")
        sys.exit(1)

    data = compile_res.json()
    output_files = data.get("outputFiles", [])

    log_entry = next((f for f in output_files if f.get("path") == "output.log"), None)
    if not log_entry:
        sys.stderr.write("No output.log found in compile artifacts.\n")
        sys.stderr.write(f"Compile status: {data.get('status')}\n")
        sys.exit(1)

    params = {
        "clsiserverid": data.get("clsiServerId", FallBack_CLSIserver),
        "compileGroup": data.get("compileGroup", "priority"),
        "editorId": payload["editorId"],
    }

    cdn_domain = data.get(
        "pdfDownloadDomain", "https://compiles.overleafusercontent.com"
    ).rstrip("/")

    # Primary URL directly to Overleaf usercontent CDN
    cdn_url = f"{cdn_domain}{log_entry['url']}"
    res = session.get(cdn_url, params=params)
    if res.status_code == 200:
        return res.text

    # Fallback to main host
    main_url = f"https://www.overleaf.com{log_entry['url']}"
    res = session.get(main_url, params=params)
    if res.status_code == 200:
        return res.text

    sys.stderr.write(f"Failed to fetch output.log (HTTP {res.status_code}).\n")
    sys.exit(1)


def audit_placeholders(log_content: str):
    has_marker = PLACEHOLDER_MARKER in log_content
    has_banner = TARGET_BANNER in log_content

    if not (has_marker or has_banner):
        sys.stdout.write("✅ Clean build: No unresolved placeholder references.\n")
        sys.exit(0)

    wrapped_re = re.compile(
        r"Package biblatex Warning:\s*Unresolved placeholder reference\s*'(?P<key>[^']+)':\s*(?P<note>[^\n]+(?:\n[^\n]+)*?\.)(?:\s+(?:Underfull|Overfull|Package|\n|\[|$))"
    )
    matches = list(wrapped_re.finditer(log_content))

    detected = {}
    if matches:
        for w in matches:
            key = w.group("key")
            note = " ".join(w.group("note").split()).rstrip(".")
            detected.setdefault(key, note)
    else:
        for line in log_content.splitlines():
            if PLACEHOLDER_MARKER in line:
                m = re.search(r"reference\s*'([^']+)':\s*(.*?)(?:\.|$)", line)
                if m:
                    detected.setdefault(m.group(1), m.group(2).strip().rstrip("."))

    total_occurrences = len(matches) if matches else len(detected)
    sys.stdout.write(
        f"\n❌ UNRESOLVED PLACEHOLDERS DETECTED ({len(detected)} unique keys, {total_occurrences} total citations):\n"
    )
    for key, note in sorted(detected.items()):
        sys.stdout.write(f"  * [{key}] {note}\n")

    sys.stdout.write("\n")
    sys.exit(2)


def main():
    project_id = os.environ.get("OVERLEAF_PROJECT_ID") or FallBack_Project_ID
    target_arg = sys.argv[1].strip() if len(sys.argv) > 1 else None

    session, csrf_token = get_authenticated_session(project_id)
    known_paths = fetch_project_paths(session, project_id, csrf_token)
    doc_mapping = fetch_project_doc_mapping(session, project_id, csrf_token)

    root_id, resource_path = resolve_file_target(target_arg, doc_mapping, known_paths)

    if target_arg:
        sys.stdout.write(f"Target Resolved:\n  * rootDoc_id: {root_id}\n  * rootResourcePath: {resource_path}\n\n")

    log_text = fetch_latest_log(
        session=session,
        project_id=project_id,
        csrf_token=csrf_token,
        root_doc_id=root_id,
        root_resource_path=resource_path,
    )
    audit_placeholders(log_text)


if __name__ == "__main__":
    main()
\end{lstlisting}


\begin{lstlisting}[language={bash}, caption={Check for remaining placholders},label=lst:bashForRemainingPlaceholders]
./fetch_overleaf_check_for_unresolved_placeholders4.py README_programmer_notes_Biber_and_BibLaTeX.tex
Target Resolved:
  * rootDoc_id: None
  * rootResourcePath: README_notes/README_programmer_notes_Biber_and_BibLaTeX.tex


❌ UNRESOLVED PLACEHOLDERS DETECTED (2 unique keys, 5 total citations):
  * [doe2025collaborative] Placeholder fake reference for use in examples
  * [maguire2026pipeline] Placeholder fake reference for use in examples
\end{lstlisting}

\item\label{itm:BuildOnGitHub} Inverting the Flow (Build on GitHub Actions via Git Bridge)

If you already maintain the two-way \textbf{Overleaf $\longleftrightarrow$ GitHub} sync, there is an alternative that avoids browser reverse-engineering altogether:
\begin{enumerate}
    \item When the student clicks \enquote{Sync to GitHub} or pushes to GitHub from Overleaf, GitHub receives the raw source tree (\enquote{.tex}, \enquote{.bib}, graphics, \etc).
    
    \item Rather than extract the log from Overleaf’s proprietary runner, let a \textbf{GitHub Actions workflow compile the document natively} using the official TeX Live / LuaLaTeX container (\enquote{texlive/texlive:latest}).
    
    \item The GitHub runner has full, unrestricted access to \enquote{stdout} and \enquote{output.log}.
    
    \item It executes the \texttt{grep} check on the spot (see \Cref{lst:bashForPlaceholders}), fails the Action check (putting a red ❌ right on the commit or pull request), and optionally drops an automated comment on the GitHub commit pointing to the unresolved keys.
\end{enumerate}

This keeps Overleaf purely as an editor, while GitHub Actions serves as the canonical verification gate.
\end{enumerate}

\begin{lstlisting}[language={bash}, caption={Check for remaining placholders},label=lst:bashForPlaceholders]
if grep -Fq "UNRESOLVED PLACEHOLDER REFERENCES DETECTED" build.log; then
    echo "Build failed: Unresolved placeholder references remain."
    exit 1
fi
\end{lstlisting}


\subsection{Further Automation}
\label{sec:furtherAutomation}
It is possible to build other automation on the mechanism shown in \Cref{lst:accessBuildLogDirectly}. For example, Python code shown in \Cref{lst:uploadOutputPDF} can copy files from the build container to the Overleaf project (\ie making them stable files in the project's directory tree). An example of using this is shown in \Cref{lst:uploadOutputSuccess}. Further examples include \fname{fordiva.json}, \fname{referencesUsed.bib}, and \fname{citedtags.bib} for a thesis.

\begin{lstlisting}[language={Python}, escapechar=£, caption={Use Python to compile the project and upload a file, such as the PDF. Note that \LaTeX{} commands were added for layout purposes.},label=lst:uploadOutputPDF]
#!/usr/bin/env python3
"""
copy_files_from_build_to_project.py

Compiles an Overleaf sub-document or project root, downloads generated build
artifacts (e.g. output.pdf, fordiva.json, referencesUsed.bib, citedtags.bib)
from the CLSI build container, and persists them into the project tree.
"""

import io
import os
import re
import sys
import time
import uuid
import requests

try:
    import browser_cookie3
except ImportError:
    browser_cookie3 = None

FallBack_Project_ID="000000000000000000000000" # replace with your project ID
FallBack_CLSIserver="clsi-cache-zone-b-prod-3" # replace with your CLSI server


def get_authenticated_session(project_id: str) -> tuple[requests.Session, str]:
    session = requests.Session()
    session.headers.update({
        "User-Agent": (
            "Mozilla/5.0 (X11; Linux x86_64) AppleWebKit/537.36 "
            "(KHTML, like Gecko) Chrome/154.0.0.0 Safari/537.36"
        ),
        "Referer": f"https://www.overleaf.com/project/{project_id}",
        "Origin": "https://www.overleaf.com",
        "Accept": "application/json, text/plain, */*",
    })

    cookie_val = os.environ.get("OVERLEAF_SESSION_COOKIE")
    if not cookie_val and browser_cookie3:
        try:
            cj = browser_cookie3.chrome(domain_name=".overleaf.com")
            session.cookies.update(cj)
            for c in cj:
                if c.name == "overleaf_session2":
                    cookie_val = c.value
                    break
        except Exception as e:
            sys.stderr.write(f"Note: browser_cookie3: {e}\n")

    if cookie_val:
        if "overleaf_session2=" in cookie_val:
            cookie_val = cookie_val.split("overleaf_session2=")[1].split(";")[0].strip()
        session.cookies.set("overleaf_session2", cookie_val, domain=".overleaf.com", path="/")

    res = session.get(f"https://www.overleaf.com/project/{project_id}")
    if res.status_code == 403 or "login" in res.url:
        sys.stderr.write("Authentication failed. Ensure you are logged into Overleaf in Chrome.\n")
        sys.exit(1)

    csrf_token = None
    for pat in [
        r'<meta\s+name=["\']ol-csrfToken["\']\s+content=["\']([^"\']+)["\']',
        r'window\.csrfToken\s*=\s*["\']([^"\']+)["\']',
        r'"_csrf"\s*:\s*["\']([^"\']+)["\']',
    ]:
        m = re.search(pat, res.text)
        if m:
            csrf_token = m.group(1)
            break

    if not csrf_token:
        sys.stderr.write("Failed to retrieve CSRF token from project page.\n")
        sys.exit(1)

    return session, csrf_token


def fetch_project_paths(session: requests.Session, project_id: str, csrf_token: str) -> list[str]:
    headers = {
        "X-Csrf-Token": csrf_token,
        "X-Requested-With": "XMLHttpRequest",
        "Accept": "application/json",
    }
    try:
        r = session.get(
            f"https://www.overleaf.com/project/{project_id}/entities?projectId={project_id}",
            headers=headers,
            timeout=10,
        )
        if r.status_code == 200:
            data = r.json()
            return [e.get("path", "").lstrip("/") for e in data.get("entities", []) if "path" in e]
    except Exception:
        pass
    return []


def get_or_create_destination_folder(
    session: requests.Session,
    project_id: str,
    csrf_token: str,
    base_folder_name: str = "artifacts",
) -> tuple[str, str]:
    """
    Creates a destination folder on the project root and returns (folder_name, folder_id).
    If base_folder_name already exists, falls back to a timestamped folder name
    to guarantee a freshly allocated, valid MongoDB ObjectId.
    """
    headers = {
        "X-Csrf-Token": csrf_token,
        "Accept": "application/json",
        "Content-Type": "application/json",
    }

    candidates = [
        base_folder_name,
        f"{base_folder_name}_{int(time.time())}",
    ]

    for fname in candidates:
        post_res = session.post(
            f"https://www.overleaf.com/project/{project_id}/folder",
            headers=headers,
            json={"name": fname},
            timeout=10,
        )
        if post_res.status_code in (200, 201):
            data = post_res.json()
            fid = data.get("_id")
            if fid:
                return fname, fid

    sys.stderr.write(f"Error: Unable to create destination folder in Overleaf.\n")
    sys.exit(1)


def compile_and_get_artifacts(
    session: requests.Session,
    project_id: str,
    csrf_token: str,
    root_resource_path: str | None = None,
) -> tuple[str, list[dict], dict]:
    compile_url = f"https://www.overleaf.com/project/{project_id}/compile?enable_pdf_caching=true"
    headers = {
        "Content-Type": "application/json",
        "X-Csrf-Token": csrf_token,
    }

    editor_id = str(uuid.uuid4())
    payload = {
        "draft": False,
        "png2pdf": True,
        "incrementalCompilesEnabled": True,
        "stopOnFirstError": False,
        "editorId": editor_id,
    }

    if root_resource_path:
        payload["rootResourcePath"] = root_resource_path
    else:
        payload["check"] = "silent"

    compile_res = session.post(compile_url, headers=headers, json=payload)
    if compile_res.status_code != 200:
        sys.stderr.write(f"POST compile failed: HTTP {compile_res.status_code}\n")
        sys.stderr.write(compile_res.text[:300] + "\n")
        sys.exit(1)

    data = compile_res.json()
    status = data.get("status", "unknown")
    routing = {
        "clsiServerId": data.get("clsiServerId", FallBack_CLSIserver),
        "compileGroup": data.get("compileGroup", "priority"),
        "editorId": editor_id,
        "pdfDownloadDomain": data.get("pdfDownloadDomain", "https://compiles.overleafusercontent.com").rstrip("/"),
    }
    return status, data.get("outputFiles", []), routing


def download_build_artifact(
    session: requests.Session,
    artifact_url: str,
    routing: dict,
) -> bytes | None:
    params = {
        "clsiserverid": routing["clsiServerId"],
        "compileGroup": routing["compileGroup"],
        "editorId": routing["editorId"],
    }

    cdn_url = f"{routing['pdfDownloadDomain']}{artifact_url}"
    res = session.get(cdn_url, params=params)
    if res.status_code == 200:
        return res.content

    main_url = f"https://www.overleaf.com{artifact_url}"
    res = session.get(main_url, params=params)
    if res.status_code == 200:
        return res.content

    return None


def upload_to_folder(
    session: requests.Session,
    project_id: str,
    csrf_token: str,
    file_bytes: bytes,
    dest_name: str,
    folder_id: str,
) -> bool:
    upload_url = f"https://www.overleaf.com/project/{project_id}/upload"
    
    # 1. folder_id MUST be in query params, not body
    params = {
        "folder_id": folder_id,
    }
    
    headers = {
        "X-Csrf-Token": csrf_token,
        "Accept": "application/json",
    }
    
    # 2. Body must contain ONLY 'name' (FineUploader uses this, validator allows this)
    data = {
        "name": dest_name,
    }
    
    # 3. File tuple: (filename, stream, mime_type)
    files = {
        "qqfile": (dest_name, io.BytesIO(file_bytes), "application/octet-stream"),
    }

    res = session.post(upload_url, headers=headers, params=params, data=data, files=files)
    if res.status_code in (200, 201):
        try:
            return res.json().get("success", True)
        except Exception:
            return True

    sys.stderr.write(f"Failed to upload '{dest_name}': HTTP {res.status_code} - {res.text[:300]}\n")
    return False


def main():
    if len(sys.argv) < 2:
        sys.stderr.write(
            "Usage:\n"
            "  ./copy_files_from_build_to_project.py [subdoc.tex] <artifact1> [artifact2 ...]\n"
            "  ./copy_files_from_build_to_project.py README_programmer_notes_Biber_and_BibLaTeX.tex output.pdf fordiva.json referencesUsed.bib citedtags.bib\n\n"
        )
        sys.exit(1)

    project_id = os.environ.get("OVERLEAF_PROJECT_ID") or FallBack_Project_ID

    args = sys.argv[1:]
    target_subdoc = None

    if args[0].endswith(".tex"):
        target_subdoc = args[0]
        targets = args[1:]
    else:
        targets = args

    if not targets:
        targets = ["output.pdf", "fordiva.json", "referencesUsed.bib", "citedtags.bib"]

    session, csrf_token = get_authenticated_session(project_id)
    known_paths = fetch_project_paths(session, project_id, csrf_token)

    # Dynamically resolve destination folder and ID
    folder_name, folder_id = get_or_create_destination_folder(session, project_id, csrf_token, "artifacts")

    root_resource_path = None
    if target_subdoc:
        clean_target = target_subdoc.strip().lstrip("./")
        root_resource_path = next(
            (p for p in known_paths if p == clean_target or os.path.basename(p) == os.path.basename(clean_target)),
            clean_target,
        )

    sys.stdout.write(f"Initiating compilation on Overleaf (target: {root_resource_path or 'default root'})...\n")
    status, artifacts, routing = compile_and_get_artifacts(
        session=session,
        project_id=project_id,
        csrf_token=csrf_token,
        root_resource_path=root_resource_path,
    )

    sys.stdout.write(f"Compilation Status: {status}\n")
    artifact_map = {item.get("path"): item.get("url") for item in artifacts if "path" in item and "url" in item}
    sys.stdout.write(f"Total build artifacts returned: {len(artifact_map)}\n")
    sys.stdout.write(f"Destination folder '{folder_name}' ID: {folder_id}\n\n")

    persisted_count = 0
    for target in targets:
        src_name = os.path.basename(target)
        dest_name = src_name

        matched_artifact_path = None
        if src_name in artifact_map:
            matched_artifact_path = src_name
        else:
            for art_path in artifact_map:
                if os.path.basename(art_path) == src_name:
                    matched_artifact_path = art_path
                    break

        if not matched_artifact_path and src_name.endswith(".pdf"):
            for art_path in artifact_map:
                if art_path.endswith(".pdf"):
                    matched_artifact_path = art_path
                    break

        if not matched_artifact_path:
            sys.stdout.write(f"⚠️  Artifact '{src_name}' was not produced by this build run (skipping).\n")
            continue

        url = artifact_map[matched_artifact_path]
        content = download_build_artifact(session, url, routing)

        if content is None:
            sys.stderr.write(f"❌ Failed to download artifact data for '{matched_artifact_path}'.\n")
            continue

        sys.stdout.write(f"£{\NotoColorEmojiFont \Uchar"1F4E6}£ Fetched '{matched_artifact_path}' ({len(content)} bytes) from build container.\n")
        success = upload_to_folder(
            session=session,
            project_id=project_id,
            csrf_token=csrf_token,
            file_bytes=content,
            dest_name=dest_name,
            folder_id=folder_id,
        )

        if success:
            sys.stdout.write(f"✅ Successfully persisted '{folder_name}/{dest_name}' into the project tree.\n")
            persisted_count += 1

    sys.stdout.write(f"\nDone. {persisted_count}/{len(targets)} files successfully persisted to Overleaf project.\n")


if __name__ == "__main__":
    main()

\end{lstlisting}

\begin{lstlisting}[language={bash}, escapechar=£, caption={Compile and upload the output.pdf into a folder in the Overleaf project. Note that \LaTeX{} commands were added for layout purposes.},label=lst:uploadOutputSuccess]
./copy_files_from_build_to_project.py README_programmer_notes_Biber_and_BibLaTeX.tex output.pdf
Initiating compilation on Overleaf (target: README_notes/README_programmer_notes_Biber_and_BibLaTeX.tex)...
Compilation Status: success
Total build artifacts returned: 177
Destination folder 'artifacts_1791413684' ID: £\hbox{6ac6cdb4a98c95d1223e6ce3}£

£{\NotoColorEmojiFont \Uchar"1F4E6}£ Fetched 'output.pdf' (311800 bytes) from build container.
£\hbox{✅ Successfully persisted}£ 'artifacts_1791413684/output.pdf' into the project tree.

Done. 1/1 files successfully persisted to Overleaf project.
\end{lstlisting}

\subsection{Risks of this automation approach}
The automation scripts introduce several technical, architectural, and operational vulnerabilities that are important to consider:
\begin{enumerate}
    \item Fragility of Undocumented Internal \glspl{API}

The scripts rely entirely on Overleaf’s private, undocumented frontend endpoints (\texttt{/updates}, \texttt{/entities}, \texttt{/compile}, \texttt{/upload}) rather than a public, supported \gls{API}. Because these endpoints power Overleaf's internal web application \gls{UI} rather than external developer integrations, \textbf{any minor frontend update by Overleaf can instantly break the scripts}. Changes to JSON schemas, \gls{URL} routing, or header requirements will cause silent failures or \texttt{HTTP 400/422} errors without warning.

    \item Local Cookie Scrape Dependency (\texttt{browser\_cookie3})

The authentication mechanism depends on reading live browser cookies (\texttt{overleaf\_session2}) from the local machine using libraries like \texttt{browser\_cookie3}. This introduces severe environment-specific constraints:
\begin{itemize}
    \item \textbf{Operating System Keychains}: Modern operating systems encrypt browser cookie databases using OS-level keyrings (macOS Keychain, Linux GNOME Keyring/KWallet), which frequently block automated Python scripts from reading them without explicit user intervention.
    
    \item \textbf{CI/CD Incompatibility}: This method \textbf{cannot} be run inside standard headless CI/CD runners (like GitHub Actions) because there is no local browser profile or active session cookie available in a clean cloud container.
\end{itemize}

    \item Hardcoded Fallbacks and Security Risks

If a user runs the script without explicitly setting the \texttt{OVERLEAF\_PROJECT\_ID} environment variable, the script defaults to interacting with a project that is statically configured in the code. While it requires authentication cookies, hardcoding project IDs or server strings (such as \texttt{clsi-cache-zone-b-prod-3}) introduces maintenance debt and a risk of cross-project interference.

    \item High Friction and Steep Learning Curve for Users

For a typical doctoral student who just wants to write a thesis, maintaining a Python script that reverse-engineers a \gls{SaaS} platform's session management creates an enormous cognitive and technical barrier. If a student changes their default browser, updates their operating system, or clears their cookies, the automation tool breaks, shifting troubleshooting overhead onto the user.

    \item The Maintenance Loop of Reverse Engineering

As seen in the script's code, handling edge cases requires complex fallback parsing (such as \texttt{regex} scans over raw text payloads, handling \texttt{FineUploader} parameters for file uploads, and timestamping fallback folders like \texttt{base\_folder\_name\_{int(time.time())}}). Maintaining this wrapper code requires continuous developer attention whenever Overleaf modifies its backend infrastructure.
\end{enumerate}

The primary vulnerability of this approach is that \textbf{it couples the thesis workflow to a closed third-party platform's private implementation details}. While it serves as a tactical workaround for Overleaf's limitations, its long-term viability depends on how long Overleaf's current internal routing remains unchanged. This underscores why \ref{itm:BuildOnGitHub} (moving the source of truth to a GitHub repository and compiling natively via GitHub Actions) shown in \Cref{sec:AutomatingPlaceholdCheck} remains the more robust, future-proof architectural choice for continuous integration.

\clearpage
\ifinswedish
\printglossary[style=mylong, type=\acronymtype, title={Lista över akronymer
och förkortningar}]
\else
%\printglossary[style=mylong, type=\acronymtype, title={List of Acronyms and abbreviations}]
%\printnoidxglossaries% Replace \printglossary[type=\acronymtype, ...] with:
\printnoidxglossary[type=\acronymtype, style=mylong, title={List of Acronyms and Abbreviations}]
\fi
\clearpage

% --- print the bibliography ---
% Print the bibliography (and make it appear in the table of contents)

% Specify the title of the references page
\renewcommand{\bibname}{References}

%\typeout{Biblatex current language is \currentlang}
\printbibliography[heading=bibintoc, title={References}]


\end{document}

